\documentclass[11pt, a4paper]{article}

\usepackage[utf8]{inputenc}
\usepackage[T1]{fontenc}
\usepackage{amsmath, amssymb, amsthm}
\usepackage{geometry}
\usepackage{hyperref}
\usepackage{xcolor}
\usepackage{braket}
\usepackage{graphicx}

\geometry{margin=1in}

\newtheorem{theorem}{Theorem}
\newtheorem{proposition}{Proposition}
\newtheorem{definition}{Definition}
\newtheorem{remark}{Remark}

\newcommand{\Hgen}{\hat{\mathcal{H}}}
\newcommand{\Kry}{\mathcal{K}}
\newcommand{\Rset}{\mathbb{R}}
\newcommand{\Cset}{\mathbb{C}}
\newcommand{\Eblr}{\mathbb{E}}
\newcommand{\Kop}{\hat{K}}

\title{\textbf{Wavefunction Flows in the \texttt{unfer} Kernel:\\
An Analytical, Grid-Free Realization of the Continuity Hamiltonian
via Orthogonal-Fock Encoding and Inversion-Free Krylov Reduction}\\[4pt]
\large A Self-Contained, Pedagogical Account}
\author{Implementation Architecture Report}
\date{July 2026 (rev.\ 5, with Tomographic Subspace Recovery, rev 14 SIRK-backed unitary flow, rev 15 module demo + benchmarks, rev 16 Quantum Bayesian Update implementation, rev 17 full SIRK-whitened Krylov basis + kernel surface wiring + 5th module demo, rev 18 productionization: P7 (README, CUDA-on-CI, MNIST 8x8 real-data validation, runtime krylov\_dim$\ge$K$_2$ check, HMC validation, 3 new bench groups) + P8.10 (CUDA kernel bench) + P8.11 (iterated\_bayes\_module 6th demo) + P9.12 (CUDA toolkit pinning) + P9.13 (CI PAT docs), rev 19 P8.9 Karcher-mean posterior estimate, rev 20--21 P8.7 Typst-math compiler + P8.8 \texttt{uk\_belief\_propagation} FFI symbol + velysterm merge: upstream typst 0.15 + KanvaSink + kanva\_typst + P9.15 partial (deletion of stale velyst examples), rev 22 P9.14 mathed\_mini Step 4, rev 23 P9.15.1 ported velyst examples, rev 24 P10.16.1 CIFAR-10 16$\times$16 baseline + SIRK shift overflow fix, rev 25 P10.16.2 real CIFAR-10 fixture + P11.24 Zenodo-Loro persistence module, rev 26 P10.16.3 rank-truncation extension (SVD on the $W$ basis, enables $d=1024$), rev 27 P10.16.4 2D Ising cross-domain validation + P10.18 $l=3$ Yang--Mills mass-gap demo + P11.22 \texttt{unfer\_edge} proxy, rev 28 P11.19/20/21/23 external-module integrations, rev 29 P10.16.5 2D Kolmogorov CFD flow-field cross-domain validation + P10.16.6 post-Newtonian gravitational-waveform cross-domain validation, rev 30 P10.16.7 $d=1024$ cross-domain scaling of all three rev 29 fixtures on the P10.16.3 rank-truncation path, rev 31 exact-projector consolidation: the $\mathcal{O}(\varepsilon)$-truncated off-diagonal coupling is removed and the exact rank-1 Mehler projector (applied via the \texttt{Operator::ProjectOnto} shortcut) becomes the only off-diagonal generator, including inside the TSR flow Hamiltonian)}

\begin{document}

\maketitle

\begin{abstract}
Layden, Sweke, Havl\'i\v{c}ek, Chowdhury, and Neklyudov~\cite{Layden2025}
recently showed that classical \emph{flow models} map exactly onto the
Schr\"odinger equation: the qsample $\Psi_t(x)=\sqrt{p_t(x)}$ of a flow's
marginal density obeys $i\,\partial_t\Psi_t=\Hgen_t\Psi_t$ for an unusual
\emph{continuity Hamiltonian} $\Hgen_t=\tfrac12[\hat p\cdot v_t(\hat x)+
v_t(\hat x)\cdot\hat p]$, which for a conservative field $v_t=\nabla V_t$
takes the elegant commutator form $\Hgen_t^c=i[\Kop,V_t(\hat x)]$. Crucially,
this generator is \emph{fully specified by the already-learned classical
velocity field} --- no further training, parameterization, or variational
optimization. Their algorithm prepares qsamples on a \emph{digital quantum
computer} by simulating $\Hgen_t^c$ with Fourier collocation and product
formulas, enabling, e.g., quadratically faster mean estimation. This report
gives a self-contained, pedagogical account of how the \texttt{unfer}
probability kernel realizes the \emph{same} wavefunction flow
\emph{classically}, with no spatial grid. We replace the $N^d$ collocation
grid by an \emph{analytical orthogonal-Fock encoding}: each data point becomes
a mutually orthogonal single-particle mode, so the conservative potential
\emph{decouples} and its flow-matching parameters reduce to a diagonal,
closed-form $\mathcal{O}(M)$ fit for $M$ data points. The resulting bounded
generator is compressed by the Shift-Invert Rational Krylov (SIRK) method of
Hashimoto and Nodera~\cite{Hashimoto2019} --- replacing the paper's product
formulas --- giving exponential approximation-error decay and reducing
single-step inference to an $\mathcal{O}(m^2)$ matrix exponential with
$m\ll M$. We also make explicit the measure-theoretic foundation: in the
infinite-dimensional limit the uniform measure on a hypersphere becomes the
Gaussian measure, and the Gegenbauer polynomials that diagonalize it become
the Hermite polynomials (L\'opez and Temme, 1999; Mehler, 1866), which
identifies the Gaussian source $p_0$ with the Fock vacuum and any orthogonal
excitation with acquired information. A second revision (June 2026) added
two further contributions. First, the exact off-diagonal generator: the Mehler vacuum,
being the uniform measure over inner wave-functions while each data channel
localizes only finitely many hyperspherical coordinates, has strictly positive
overlap $\varepsilon_j=\braket{0|x_j}$ with every channel, and that
non-orthogonality alone drives the transport. The exact off-diagonal generator is
therefore \emph{just the vacuum projector} $\Hgen_{\mathrm{off}}=\ket{0}\bra{0}$
--- Hermitian, idempotent, with closed-form evolution
$e^{-it\Hgen_{\mathrm{off}}}=I+(e^{-it}-1)\ket{0}\bra{0}$ --- realizing
coherent amplitude transport between the source and the data channels in
place of the stationary diagonal surrogate, with \emph{no}
$\mathcal{O}(\varepsilon)$ truncation: the kernel applies it through the
rank-1 shortcut $\Hgen_{\mathrm{off}}\ket{s}=\braket{0|s}\ket{0}$, never
an expanded coupling-term list.
Second, the full \emph{Tomographic Subspace Recovery} (TSR) pipeline: a
two-level Count-Sketch hashing scheme ($S_1:\Rset^d\to\Rset^k$, $S_2$:
features $\to$ $K_2$-dim single-excitation Fock states) that decouples
semantics from reasoning, pre-projects all raw-coordinate observables into
the $m^2$ Krylov operator basis, and executes online generation in four
phases (encode, evolve, tomographic reconstruct, lossless decode) with
\emph{no} $\mathcal{O}(M)$ dependence at inference. We close by detailing
the integration into \texttt{unfer}'s GPU-accelerated \texttt{fock\textunderscore
sirk} solver, its \texttt{nested\textunderscore fock\textunderscore algebra}
CAS, the new \texttt{qfm} tomographic crate, and the Austral module runtime.
Subsequent revisions (through rev 28, July 2026) close the pipeline
empirically: real-data validation on MNIST $8\times 8$ and CIFAR-10
$16\times 16$, a rank-truncation extension that removes the
$\mathcal{O}(K_2^3)$ observable-projection wall (enabling $d=1024$),
cross-domain validation on critical 2D Ising configurations, and an $l=3$
SU(2) lattice Yang--Mills mass-gap demonstration within 1\% of the
strong-coupling estimate.
\end{abstract}

\tableofcontents

\section{Introduction}

\subsection{Flow models and the qsample bridge}

\emph{Flow models} are a cornerstone of modern generative machine learning.
They learn a continuous-time Markov process --- realized on an ordinary
(classical) computer --- that transports a simple source distribution $p_0$
(typically isotropic Gaussian noise) into a complex target $p_{\mathrm{targ}}$
from which we are given samples. The process is specified by a time-dependent
\emph{velocity field} $v_t(x)$ through the ordinary differential equation (ODE)
\begin{equation}
  \frac{\mathrm{d}}{\mathrm{d}t}\,x_t=v_t(x_t),\qquad x_0\sim p_0,
  \label{eq:flowode}
\end{equation}
whose random initial condition makes $x_t\sim p_t$ a random variable. The
family of marginals $(p_t)_{t\in[0,T]}$ is the \emph{probability path}, and by
the change-of-variables formula it satisfies the \emph{continuity equation}
\begin{equation}
  \partial_t\,p_t(x)=-\nabla\!\cdot\bigl[v_t(x)\,p_t(x)\bigr].
  \label{eq:continuityeq}
\end{equation}
Flow models learn $v_t^\theta\approx v_t$ (e.g.\ by flow matching,
Section~\ref{sec:flowprimer}) so that integrating~\eqref{eq:flowode} from
$x_0\sim p_0$ yields $x_T\sim p_T\approx p_{\mathrm{targ}}$.

In quantum computing, the corresponding object of interest is the
\emph{qsample}: the coherent encoding
\begin{equation}
  \ket{\Psi_p}=\sum_x\sqrt{p(x)}\,\ket{x}
  \quad\text{(discrete)},\qquad
  \ket{\Psi_p}=\int\sqrt{p(x)}\,\ket{x}\,\mathrm{d}x
  \quad\text{(continuous)}.
\end{equation}
Measuring $\ket{\Psi_p}$ in the $\{\ket{x}\}$ basis returns a sample
$x\sim p$, but \emph{preparing} a qsample is strictly more powerful than
sampling: with qsample access one can build quantum algorithms for mean
estimation and property testing with provable sample-complexity advantages
(Section~\ref{sec:qsample}). Preparing qsamples for arbitrary $p$ is hard in
general, so the central question is: \emph{for which $p$ can we prepare
qsamples efficiently?}

\subsection{Wavefunction flows (Layden et al.)}

Layden, Sweke, Havl\'i\v{c}ek, Chowdhury, and Neklyudov~\cite{Layden2025}
answer this for the vast family of distributions expressible by flow models.
Their key observation is a striking identity: if $p_t$ obeys the continuity
equation~\eqref{eq:continuityeq}, then its qsample wavefunction
$\Psi_t(x)=\sqrt{p_t(x)}$ obeys the Schr\"odinger equation
$i\,\partial_t\Psi_t=\Hgen_t\Psi_t$ for the \emph{continuity Hamiltonian}
\begin{equation}
  \Hgen_t=\tfrac12\bigl[\hat p\cdot v_t(\hat x)+v_t(\hat x)\cdot\hat p\bigr],
  \label{eq:continuityH}
\end{equation}
derived in Section~\ref{sec:continuity}. The dynamics it generates they call a
\emph{wavefunction flow}. Two features make this remarkable:
\begin{enumerate}
\item \textbf{No further learning.} $\Hgen_t$ is \emph{completely determined by
the already-learned classical velocity field} $v_t^\theta$. There is no
parameterized quantum circuit, no variational loop, no barren-plateau
pathology --- the entire ``training'' has been done classically, and the
quantum side is pure Hamiltonian simulation.
\item \textbf{Provable qsample preparation.} Evolving the easily prepared
initial qsample $\Psi_0=\sqrt{p_0}$ (a tensor product of standard-normal
qsamples for Gaussian $p_0$) under $\Hgen_t$ for time $T$ produces exactly
$\Psi_T=\sqrt{p_T}$; the distance to the ideal target $\sqrt{p_{\mathrm{targ}}}$
is controlled \emph{entirely classically} by $\|p_T-p_{\mathrm{targ}}\|$.
\end{enumerate}
The paper's main technical work is then turning this into a digital quantum
algorithm: discretizing space (Fourier collocation on an $N^d$ grid) and time
(product formulas), with error bounds (their Theorems~1 and~2).

\subsection{This work: a classical, grid-free realization in \texttt{unfer}}

The Layden et al.\ route targets a \emph{digital quantum computer} and pays for
it with an $N^d$ spatial collocation grid and a product-formula time
discretization whose cost grows with the operator norm $\|H_t^c\|=
\mathcal{O}(N^2 d/L^2)$. The \texttt{unfer} probability kernel already carries
all the machinery to realize the \emph{same} continuity-Hamiltonian
wavefunction flow \emph{classically}, on a separable Fock space, with Born-rule
semantics. This report shows how, and in doing so removes the spatial grid
altogether. The construction rests on three pillars, developed pedagogically
in Sections~\ref{sec:continuity}--\ref{sec:krylov}:

\begin{enumerate}
\item \textbf{Orthogonal Fock encoding (Section~\ref{sec:framework}).} Instead
of an $N^d$ collocation grid, encode each of the $M$ data points as a distinct,
mutually orthogonal single-boson mode. For a conservative field the potential
$V_t$ then becomes \emph{diagonal}, the flow-matching normal equations decouple
into $M$ independent scalar equations, and ``training'' is an $\mathcal{O}(M)$
closed-form average with zero data loss --- grid-free and inversion-free.

\item \textbf{Inversion-free Krylov reduction (Section~\ref{sec:krylov}).} In
place of the paper's product-formula time stepping, compress the resulting
time-independent generator $\bar H$ with the Shift-Invert Rational Krylov
(SIRK) method of Hashimoto and Nodera~\cite{Hashimoto2019}, whose error decays
\emph{exponentially} in the subspace dimension $m$. Because the potential has
no cross-terms, each Krylov step costs $\mathcal{O}(M)$ and inference is a
single $\mathcal{O}(m^2)$ matrix exponential.

\item \textbf{The Gaussian-source $=$ Fock-vacuum identification
(Section~\ref{sec:mehler}).} The Mehler/L\'opez--Temme limit of Gegenbauer to
Hermite polynomials makes precise why the Gaussian source $p_0$ of a flow model
is exactly the Fock vacuum, and why orthogonal excitations carry information.
\end{enumerate}

Section~\ref{sec:impl} maps every piece onto the existing \texttt{unfer}
modules. Throughout we keep prerequisites minimal: a reader comfortable with
ladder operators, Gaussian integrals, and the matrix exponential should be able
to follow without external references. We are explicit about scope: the
shipped \texttt{unfer} builtins implement a spectrally-stable diagonal
surrogate and the \emph{exact} off-diagonal generator --- which turns out
to be just the projector $\ket{0}\bra{0}$ onto the Mehler vacuum,
off-diagonal purely through the vacuum's nonzero overlap with the
localized data channels, applied via its rank-1 structure with no
truncated coupling expansion --- and we flag the differences precisely
(Sections~\ref{sec:framework}, \ref{sec:innerouter}, and~\ref{sec:impl}).

\section{Background I: flow matching and the qsample}
\label{sec:flowprimer}

\subsection{Learning the velocity field by flow matching}

A flow model needs a velocity field $v_t$ transporting $p_0$ to
$p_{\mathrm{targ}}$. The ground-truth $v_t$ that generates a chosen probability
path is unknown, so the naive flow-matching loss is intractable. The decisive
trick~\cite{Layden2025} is the \emph{conditional flow matching} (CFM) loss,
which has the same gradient but is a tractable regression. For the
straight-line (optimal-transport) interpolation
$x_t=(1-t)x_0+t\,x_{\mathrm{targ}}$ with constant conditional velocity
$x_{\mathrm{targ}}-x_0$,
\begin{equation}
  \mathcal{L}_{\mathrm{CFM}}(\theta)=
  \Eblr_{t\sim\mathrm{unif}[0,1]}\,
  \Eblr_{x_0\sim p_0}\,
  \Eblr_{x_{\mathrm{targ}}\sim p_{\mathrm{targ}}}
  \Bigl\|v_t^\theta\bigl[(1-t)x_0+t\,x_{\mathrm{targ}}\bigr]
  -\bigl(x_{\mathrm{targ}}-x_0\bigr)\Bigr\|^2 .
  \label{eq:cfm}
\end{equation}
This estimates the gradient of the true flow-matching loss by an unbiased Monte
Carlo average over training samples and \emph{never} solves a differential
equation. The wavefunction-flow construction is agnostic to \emph{how} $v_t$ is
learned; it only requires that some $v_t^\theta\approx v_t$ exists. In
Section~\ref{sec:framework} we exploit~\eqref{eq:cfm} directly: under the
orthogonal-Fock encoding its minimizer is available in closed form.

\subsection{Diffusion models are flow models}

Diffusion models~\cite{Layden2025} provide another route to $v_t^\theta$. A
forward stochastic differential equation $\mathrm{d}x_t=f_t(x_t)\,\mathrm{d}t+
g_t\,\mathrm{d}w_t$ drives $p_{\mathrm{targ}}$ to a Gaussian; the associated
\emph{probability-flow ODE} replaces the reverse SDE by a deterministic field
\begin{equation}
  v_t(x)=f_t(x)-\tfrac12 g_t^2\,\nabla\ln p_t(x),
\end{equation}
which is an instance of~\eqref{eq:flowode}. When $f_t$ is conservative so is
$v_t$, so diffusion models fall under the same conservative-field analysis we
use below. This is why the single construction covers both of the dominant
continuous-variable generative paradigms.

\subsection{Qsamples and their statistical payoff}
\label{sec:qsample}

The reason to lift $p$ to its qsample $\Psi=\sqrt p$ is statistical leverage.
For mean estimation of $f(x)$ under $x\sim p$, the best classical estimators
(sub-Gaussian, e.g.\ median-of-means) achieve, with $m$ samples and failure
probability $\Delta$,
\begin{equation}
  \Pr\Bigl[|\tilde\mu-\mu|>C\sigma\sqrt{\tfrac{\log(1/\Delta)}{m}}\Bigr]\le\Delta,
  \qquad \sigma=\sqrt{\mathrm{Var}_{x\sim p}f(x)} .
\end{equation}
Given qsample-preparation access, the quantum sub-Gaussian
estimator~\cite{Layden2025} instead achieves
\begin{equation}
  \Pr\Bigl[|\tilde\mu-\mu|>\sigma\,\tfrac{\log(1/\Delta)}{m}\Bigr]\le\Delta,
\end{equation}
a \emph{quadratic} speed-up in the dependence on $m$, with no promise on the
variance beyond finiteness. Wavefunction flows make exactly this access
available for any distribution expressible by a (conservative) flow model. The
\texttt{unfer} kernel exposes the analogous probabilities through its Born-rule
layer, so the same statistical quantities are computable on the classical
substrate.

\section{Background II: operator functions and the SIRK method}
\label{sec:sirk}

Realizing a wavefunction flow --- on a quantum computer or on \texttt{unfer} ---
ultimately asks us to apply $e^{-i\Hgen t}$ (or, more generally, an operator
$\varphi$-function) to a state. The paper~\cite{Layden2025} does this with
product formulas; \texttt{unfer} does it with the Shift-Invert Rational Krylov
(SIRK) method of Hashimoto and Nodera~\cite{Hashimoto2019}, which we summarize
here because it is the engine of Section~\ref{sec:krylov}.

\subsection{Exponential integrators and \texorpdfstring{$\varphi$}{phi}-functions}

Evolution equations $u'=Au+g$ have the variation-of-constants solution
\begin{equation}
  u(t)=e^{tA}u_0+\int_0^t e^{(t-s)A}g(s)\,\mathrm{d}s .
\end{equation}
\emph{Exponential integrators} discretize the integral using the family of
$\varphi$-functions
\begin{equation}
  \varphi_0(z)=e^z,\qquad
  \varphi_k(z)=\int_0^1 e^{sz}\,\frac{(1-s)^{k-1}}{(k-1)!}\,\mathrm{d}s
  \quad(k\ge 1),
\end{equation}
e.g.\ the one-step scheme $u_i=e^{\delta A}u_{i-1}+\delta\,\varphi_1(\delta
A)\,g_{i-1}$. The computational kernel is therefore the product
$\varphi_k(A)\,v$ of an operator function with a vector --- exactly the
primitive that wavefunction-flow inference needs (with $k=0$, $A=-i\Hgen$).

\subsection{Why naive (Arnoldi) Krylov fails for stiff generators}

The Arnoldi method builds the polynomial Krylov subspace
\begin{equation}
  \Kry_m(A,v)=\mathrm{span}\{v,Av,\dots,A^{m-1}v\},
\end{equation}
and approximates $\varphi_k(A)v\approx V_m\,\varphi_k(\mathbf{H}_m)\,V_m^*v$,
where $V_m$ holds an orthonormal basis and
$\mathbf{H}_m=V_m^*AV_m$ is the compression of $A$. For a \emph{bounded} $A$
the error decays geometrically, governed by polynomial approximation of
$\varphi_k$ on the numerical range $W(A)$. But the continuity Hamiltonian
contains the kinetic operator $\Kop=-\tfrac12\nabla^2$, which is
\emph{unbounded}: $W(A)$ is unbounded, the relevant conformal map ceases to
exist, and the error is no longer guaranteed to decay in $m$. This is the
stiffness obstruction --- the same one the paper~\cite{Layden2025} confronts
when it notes $\|K\|=\mathcal{O}(N^2 d/L^2)$ blows up with grid resolution.

\subsection{Resolvents, rational Krylov, and SIRK}

The cure is to work with the \emph{resolvent} $X=(\gamma I-A)^{-1}$, which is
bounded whenever $\gamma$ is outside the spectrum of $A$. One can define
operator functions of an unbounded $A$ through the bounded $X$ via
$f(A):=f_\gamma(X)$ with $f_\gamma(z)=f(\gamma-z^{-1})$. The
\emph{Rational Krylov} subspace uses a sequence of shifts $\gamma_j$ and
resolvents $X_j=(\gamma_j I-A)^{-1}$,
\begin{equation}
  \mathcal{Q}_m(\{X_j\},v)
  =\mathrm{span}\{v,\,X_1v,\,X_2X_1v,\dots,X_{m-1}\cdots X_1 v\}.
  \label{eq:ratkrylov}
\end{equation}
The Shift-Invert Rational Krylov (SIRK) method chooses the real, equispaced
shifts $\gamma_j=N-hj$ with $N>hj>0$. A key algebraic identity is that
all the resolvents share a single ``base'' resolvent $X_m$,
\begin{equation}
  X_j=(\gamma_j I-A)^{-1}=\bigl(I+h(m-j)X_m\bigr)^{-1}X_m,
\end{equation}
so that the rational Krylov space~\eqref{eq:ratkrylov} equals
$\{r(X_m)v:\,r\in\mathcal{R}^{\mathrm{SIRK}}_{m-1}\}$ with denominators
$q_m(z)=(1+hmz)\cdots(1+hz)$. Hashimoto and Nodera prove an error bound that,
crucially, decays \emph{exponentially}:
\begin{equation}
  \bigl\|\varphi_k(A)v-V_m\,\psi_{k,\gamma_m}(\mathbf{H}_m\mathbf{K}_m^{-1})\,V_m^*v\bigr\|
  \;\le\;2C\,\|v\|\,e^{-hm}\!\!\min_{r\in\mathcal{R}^{\mathrm{SIRK}}_{m-1}}\!
  \|f_{k,N}-r\|_{\infty,\Sigma},
  \label{eq:sirkbound}
\end{equation}
with $C\in[2,11.08]$ and $\Sigma$ a bounded convex hull of resolvent numerical
ranges. The factor $e^{-hm}$ is the practical payoff: where Arnoldi (for
unbounded $A$) and even plain rational Krylov decay only polynomially, SIRK
decays geometrically, so a small subspace dimension $m$ (e.g.\ $m=15$--$20$)
suffices. This is the engine \texttt{unfer} already ships as
\texttt{fock\textunderscore sirk}; Section~\ref{sec:krylov} explains why our
decoupled generator makes it run especially fast.

\section{The continuity Hamiltonian}
\label{sec:continuity}

We now derive, following Layden et al.~\cite{Layden2025}, the generator at the
heart of a wavefunction flow, and record the conservative special case that
\texttt{unfer} implements.

\subsection{General case}

Let $\Psi_t(x)=\sqrt{p_t(x)}$ be the qsample wavefunction of a probability path
satisfying the continuity equation~\eqref{eq:continuityeq}. Differentiating and
substituting,
\begin{align}
  i\,\partial_t\Psi_t(x)
  &=\frac{i}{2\sqrt{p_t(x)}}\,\partial_t p_t(x)
  =-\frac{i}{2\Psi_t(x)}\,\nabla\!\cdot\bigl[v_t(x)\,\Psi_t(x)^2\bigr]\notag\\
  &=\tfrac12\Bigl\{(-i\nabla)\cdot\bigl[v_t(x)\Psi_t(x)\bigr]
  +v_t(x)\cdot(-i\nabla)\Psi_t(x)\Bigr\}
  =\bigl(\Hgen_t\Psi_t\bigr)(x).
  \label{eq:Hderiv}
\end{align}
Thus $\Psi_t=\sqrt{p_t}$ obeys $i\,\partial_t\Psi_t=\Hgen_t\Psi_t$ with the
\emph{continuity Hamiltonian}
\begin{equation}
  \Hgen_t=\tfrac12\bigl[\hat p\cdot v_t(\hat x)+v_t(\hat x)\cdot\hat p\bigr],
  \qquad
  (\hat x_j\Psi)(x)=x_j\Psi(x),\quad
  (\hat p_j\Psi)(x)=-i\,\partial_{x_j}\Psi(x).
  \label{eq:Hgeneral}
\end{equation}
The symmetrization $\tfrac12[\hat p\cdot v+v\cdot\hat p]$ is what makes
$\Hgen_t$ Hermitian even though $\hat p$ and $v_t(\hat x)$ do not commute; it is
the operator-ordering that exactly reproduces the divergence form of the
continuity equation.

\subsection{Conservative special case}

By a theorem on absolutely continuous curves in 2-Wasserstein space, every
probability path admits a \emph{conservative} velocity field
$v_t(x)=\nabla V_t(x)$ for a scalar potential $V_t$, without loss of
generality. Substituting $v_t=\nabla V_t$ into~\eqref{eq:Hderiv} and using
$\nabla^2(V_t\Psi_t)=\Psi_t\nabla^2 V_t+2\nabla V_t\cdot\nabla\Psi_t
+V_t\nabla^2\Psi_t$ gives the compact \emph{commutator form}
\begin{equation}
  \Hgen_t^c=i\,[\Kop,\,V_t(\hat x)],\qquad
  \Kop=\tfrac12\,\hat p\cdot\hat p,\quad
  (\Kop\Psi)(x)=-\tfrac12\nabla^2\Psi(x),
  \label{eq:Hconservative}
\end{equation}
which acts as
\begin{equation}
  \bigl(\Hgen_t^c\Psi_t\bigr)(x)
  =-\frac{i}{2}\Bigl[\Psi_t(x)\,\nabla^2 V_t(x)
  +2\,\nabla V_t(x)\cdot\nabla\Psi_t(x)\Bigr].
  \label{eq:Hcaction}
\end{equation}
The two terms are the source/compressibility ($\nabla^2 V_t$) and the drift
($\nabla V_t\cdot\nabla$) of the flow. Equation~\eqref{eq:Hconservative} is the
operator \texttt{unfer} targets; note the drift carries coefficient $2$
relative to the source, exactly as in~\eqref{eq:Hcaction}.

\subsection{Context: the Madelung decoupling}

It is worth contrasting $\Hgen_t$ with the standard Hamiltonian
$\Kop+V_t(\hat x)$. The Madelung equations re-express the dynamics of the
latter as a continuity equation for $|\Phi_t|^2$ \emph{coupled} to a separate
equation for the phase $\arg\Phi_t$: density and phase feed back on each other.
The continuity Hamiltonian is engineered to \emph{decouple} them --- the phase
stays fixed while the velocity field is specified directly. The price is that
$\Hgen_t$ is not of canonical kinetic-plus-potential form and so ``may not occur
in nature''~\cite{Layden2025}; the reward is that it can nonetheless be
simulated efficiently, and that $\Psi_t$ stays a bona fide qsample
($\Psi_t=\sqrt{p_t}\ge 0$, no spurious phase) at every time.

\section{Theoretical framework: orthogonal Fock encoding}
\label{sec:framework}

We now diverge from the digital-quantum route of~\cite{Layden2025} --- which
discretizes $V_t(\hat x)$ and $\Kop$ on an $N^d$ Fourier grid --- and instead
encode the data so that the conservative continuity Hamiltonian becomes
\emph{diagonal and grid-free}. This is the analytical core of the
\texttt{unfer} realization.

\subsection{Zero-loss data encoding and the symmetric Fock space}

To avoid the $\mathcal{O}(M^3)$ matrix inversion that plagues analytical kernel
methods, we treat the $M$ training points as events of a Poisson counting
measure. Via the Wiener--It\^o--Segal isomorphism, each data point
$x_j\in\Rset^D$ is mapped to a \emph{strictly orthogonal} single-particle
state of a symmetric Fock space:
\begin{equation}
  \ket{x_j}=\hat c_j^\dagger\ket{0},\qquad
  \braket{x_i | x_j}=\delta_{ij}.
\end{equation}
In the coordinate representation these correspond to localized, normalized
wave-packets $\Psi_j(x)=\braket{x | x_j}$ centered at $x_j$. Because the
events of a point process are disjoint, the packets have \emph{disjoint
supports}, which gives the two identities the whole construction exploits:
\begin{align}
  \Psi_j(x)\,\Psi_k(x)&=\delta_{jk}\,\Psi_j(x)^2,
  \label{eq:disjoint}\\
  \nabla\Psi_j(x)\cdot\nabla\Psi_k(x)&=0\quad(j\neq k).
  \label{eq:graddisjoint}
\end{align}
The orthogonality~\eqref{eq:disjoint} is ``zero data loss'': no two data points
share amplitude. The gradient identity~\eqref{eq:graddisjoint} is what
decouples the optimization. This is the grid-free counterpart of the paper's
Fourier collocation: the $N^d$ grid points are replaced by $M$ orthogonal
modes, one per datum.

\paragraph{Inner vs.\ outer Fock space, and how Mehler connects them.} The
kernel's symbolic engine, \texttt{nested\textunderscore fock\textunderscore
algebra}, implements this construction as two syntactically distinct levels
(Section~\ref{sec:innerouter} makes the connection between them precise).
\emph{Level 1, the inner Fock space}, is the Hilbert space \emph{of
wave-functions}: each of its elements is itself a wave-function/qsample over
$x\in\Rset^D$, and its normalized elements are exactly the points of the
infinite-dimensional hypersphere of Section~\ref{sec:mehler}. \emph{Level 2,
the outer Fock space} (the ``multiverse'') is built directly on top of the
inner space by Mehler's correspondence (Section~\ref{sec:innerouter}): the
Gegenbauer polynomials that diagonalize the uniform measure on that
hypersphere of inner wave-functions become, in the infinite-dimensional
limit, the Hermite polynomials $H_n$ --- and the $H_n$ are exactly the
number-state basis of a bosonic Fock space, $H_0$ the vacuum $\ket{0}$ and
each higher order a further excitation. Each packet $\Psi_j$, an element of
the inner space, thus indexes a one-particle ket
$\ket{x_j}=\hat c_j^\dagger\ket{0}$ of this outer space. This is literal in
\texttt{nested\textunderscore fock\textunderscore algebra}: an
\texttt{OuterState} is a multiset of inner-space configurations
(\texttt{InnerBosonicState}/\texttt{InnerFermionicState}), so an outer boson's
``mode label'' is an entire inner-space configuration, not a bare integer.
Every $\ket{0}$, $H_0$, and $P_0$ used operationally from here through
Section~\ref{sec:impl} --- including the $\bar H_{\text{impl}}$ and
$\Hgen_{\text{off}}$ of Section~\ref{sec:impl} --- is this \emph{outer}
vacuum. Section~\ref{sec:innerouter} shows precisely how Mehler's
Gegenbauer-to-Hermite limit constructs it, together with the whole basis,
from the uniform prior over inner wave-functions --- and computes the small
but strictly \emph{positive} overlap
$\braket{0|x_j}=\varepsilon_j$~\eqref{eq:mehleroverlap} between that
uniform vacuum and each localized data channel, the non-orthogonality that
Section~\ref{sec:impl} turns into the off-diagonal transport mechanism
itself.

\subsection{Analytical potential optimization}

Write the conservative flow potential of~\eqref{eq:Hconservative} as a linear
combination of the orthogonal packets,
\begin{equation}
  V_t(x)=\sum_{j=1}^M\alpha_j(t)\,\Psi_j(x).
\end{equation}
Fitting the conservative velocity $\nabla V_t$ to the straight-line target
$v_t=x_{\mathrm{targ}}-x_0$ is the conditional-flow-matching least-squares
problem~\eqref{eq:cfm}, which here reads
\begin{equation}
  \mathcal{L}(\vec\alpha(t))=
  \Eblr_{t,x_0,x_{\mathrm{targ}}}\Bigl\|
  \sum_{j=1}^M\alpha_j(t)\,\nabla\Psi_j(x_t)-(x_{\mathrm{targ}}-x_0)\Bigr\|_2^2 ,
  \qquad x_t=(1-t)\,x_0+t\,x_{\mathrm{targ}},
\end{equation}
where $x_{\mathrm{targ}}$ is drawn uniformly from the $M$ training points
$\{x^{(1)},\dots,x^{(M)}\}$. Setting $\partial\mathcal{L}/\partial\alpha_k=0$
gives the normal equations
\begin{equation}
  \sum_{j=1}^M\alpha_j(t)\,
  \Eblr_{x_t}\!\bigl[\nabla\Psi_j(x_t)\cdot\nabla\Psi_k(x_t)\bigr]
  =\Eblr_{x_0,x_{\mathrm{targ}}}\!\bigl[(x_{\mathrm{targ}}-x_0)\cdot
  \nabla\Psi_k(x_t)\bigr].
\end{equation}
By the gradient-disjointness~\eqref{eq:graddisjoint} the left-hand matrix is
\emph{strictly diagonal}, and on any path toward a target
$x_{\mathrm{targ}}=x^{(j)}$ with $j\neq k$ the packet gradient
$\nabla\Psi_k(x_t^{(j)})$ vanishes. The system collapses to a closed-form,
per-mode solution:
\begin{equation}
  \alpha_k(t)=\frac{\Eblr_{x_0}\!\bigl[(x^{(k)}-x_0)\cdot
  \nabla\Psi_k(x_t^{(k)})\bigr]}
  {M\,\Eblr_{x_0}\!\bigl[\|\nabla\Psi_k(x_t^{(k)})\|^2\bigr]},
  \label{eq:alpha}
\end{equation}
with $x_0\sim\mathcal{N}(0,\mathbf{I})$ the Gaussian source $p_0$ and
$x_t^{(k)}=(1-t)x_0+t\,x^{(k)}$. Each $\alpha_k$ is an independent scalar
average, so the whole ``training'' is $\mathcal{O}(M)$ with zero data loss ---
no inner optimization loop, and no $N^d$ grid to sweep. This is the concrete
payoff of the orthogonal encoding: it turns the CFM regression~\eqref{eq:cfm}
into $M$ decoupled one-dimensional integrals.

\subsection{The continuity generator and a vacuum-projector prior}

Substituting $V_t=\sum_j\alpha_j(t)\Psi_j$ into the conservative continuity
Hamiltonian~\eqref{eq:Hconservative} yields the off-diagonal generator
\begin{equation}
  \Hgen_t^c=-\frac{i}{2}\sum_{j=1}^M\alpha_j(t)\,\hat h_j(x,\nabla),\qquad
  \hat h_j=2\,\nabla\Psi_j\cdot\nabla+\nabla^2\Psi_j ,
  \label{eq:Hgen_full}
\end{equation}
in which $\hat h_j$ transports amplitude into the data channel $\ket{x_j}$.

\paragraph{The source prior, and an \texttt{unfer} modeling extension.}
In Layden et al.~\cite{Layden2025} the source $p_0$ enters \emph{only} through
the initial qsample $\Psi_0=\sqrt{p_0}$; it is not part of the Hamiltonian. In
the \texttt{unfer} realization we find it convenient to also encode a
``no-information'' prior \emph{inside} the generator, as the rank-1 projector
onto the Mehler ground state,
\begin{equation}
  H_0=\ket{0}\bra{0},
\end{equation}
which represents total prior uncertainty (made precise in
Section~\ref{sec:mehler}). This term is \emph{far from inert}: because the
Mehler vacuum has strictly positive overlap
$\varepsilon_j=\braket{0|x_j}$~\eqref{eq:mehleroverlap} with every localized
data channel, the projector $H_0$ is \emph{by itself} an off-diagonal
generator in the channel basis
($\bra{x_i}H_0\ket{x_j}=\varepsilon_i\varepsilon_j$), and
Section~\ref{sec:impl} shows it is the \emph{exact} off-diagonal
generator~\eqref{eq:Hmehler} --- the vacuum$\,\leftrightarrow\,$channel
transport needs nothing beyond it. Only in the $\varepsilon\to0$ frame
approximation (the code's \texttt{ProjectVacuum}) does $H_0$ degenerate to
an inert, self-adjoint, idempotent term whose only effect is a phase on the
vacuum. We are explicit that this projector term is an \texttt{unfer}
modeling choice, \emph{not} part of the paper's continuity Hamiltonian; it
can be set to zero to recover the pure generator~\eqref{eq:Hgen_full}.

\subsection{Time-averaging and exact commutativity}

Because the packets have disjoint supports, the localized operators act on
independent regions and \emph{exactly commute}:
\begin{equation}
  [\hat h_j,\hat h_k]=0\quad\forall j,k
  \;\Longrightarrow\;[\Hgen_t,\Hgen_{t'}]=0 .
\end{equation}
The time-ordering operator therefore drops out of the evolution, and we may
analytically time-average the parameters, $\bar\alpha_j=\int_0^1\alpha_j(t)\,
\mathrm{d}t$. The continuous-time flow collapses to a single,
\emph{time-independent} generator,
\begin{equation}
  \bar H=\ket{0}\bra{0}-\frac{i}{2}\sum_{j=1}^M\bar\alpha_j\,\hat h_j ,
  \label{eq:Hbar}
\end{equation}
and inference is one exact step, $\Psi_{t=1}=e^{-i\bar H}\Psi_0$. (Setting the
projector term to zero recovers the paper's generator exactly; we keep it for
the generalization behavior described above.) This commutativity is the
analytic analogue of the paper's observation that a conveniently smooth
probability path can be chosen to tame the simulation error.

\paragraph{The implemented (diagonal) generator.}
The differential operators $\hat h_j$ of~\eqref{eq:Hbar} mix the vacuum with
the data channels; that off-diagonal generator is the full wavefunction flow.
The \texttt{unfer} builtin (Section~\ref{sec:impl}) currently implements the
\emph{diagonal} number-operator surrogate
\begin{equation}
  \bar H_{\text{impl}}=\ket{0}\bra{0}+\sum_{j=1}^M\bar\alpha_j\,\hat n_j,
  \qquad \hat n_j=\hat c_j^\dagger\hat c_j ,
  \label{eq:Hdiag}
\end{equation}
for which the vacuum and each data channel are eigenstates
($\bar H_{\text{impl}}\ket{0}=\ket{0}$,
$\bar H_{\text{impl}}\ket{x_j}=\bar\alpha_j\ket{x_j}$) --- exactly so when
$\ket{0}$ is read as the code's frame vacuum
(\texttt{ProjectVacuum}), and to $\mathcal{O}(\varepsilon)$ for the true
Mehler vacuum of Section~\ref{sec:innerouter}. This is the
spectrally-stable skeleton of~\eqref{eq:Hbar}: it shares the prior, the
$\mathcal{O}(M)$ construction, and the SIRK compressibility, while the
off-diagonal mixing --- the part that actually performs the wavefunction
flow --- is supplied exactly by the Mehler projector itself
(eq.~\eqref{eq:Hmehler}, Section~\ref{sec:impl}), with no truncated
companion form.

\section{Relating the Fock space with a uniform measure in an
infinite-dimensional sphere}
\label{sec:mehler}

This section justifies, from first principles, the central modeling choice:
\emph{the Gaussian source $p_0$ of a flow model is the uniform
``no-information'' measure, the Fock vacuum is its ground state, and orthogonal
excitations encode information.} The bridge is a classical limit, due in spirit
to Mehler~(1866) and made quantitative by L\'opez and
Temme~\cite{LopezTemme1999}: in the infinite-dimensional limit the uniform
measure on a hypersphere becomes the Gaussian measure, and the Gegenbauer
polynomials --- which are orthogonal with respect to the uniform
(hyperspherical) measure and define the hyperspherical harmonics --- become
the Hermite polynomials. This is what makes the easily-prepared Gaussian
initial qsample $\Psi_0=\sqrt{p_0}$ of Layden et al.\ coincide with the Fock
vacuum that \texttt{unfer} manipulates natively.

\subsection{Uniform information is the hyperspherical measure}

Consider the surface of a $k$-sphere parameterized by angles
$\varphi_0,\dots,\varphi_{k-1}$ with $x_j=\cos\varphi_j$. Its rotation-invariant
(uniform) area element factorizes as
\begin{equation}
  \mathrm{d}\sigma(k)=\prod_{j=0}^{k-1}
  \bigl(\sqrt{1-x_j^2}\bigr)^{\,j-1}\,\mathrm{d}x_j ,
  \label{eq:areaelement}
\end{equation}
where the $j=0$ factor is special: $\mathrm{d}x_0=\sin\varphi_0\,
\mathrm{d}\varphi_0$ and the azimuthal angle $\varphi_0$ ranges over
$[0,2\pi]$, while every other angle ranges over $[0,\pi]$. The weight
$(1-x^2)^{(\alpha-1)/2}$ appearing in a single coordinate of
\eqref{eq:areaelement} is exactly the Gegenbauer weight of order
$\alpha$. Because the Gegenbauer polynomials are orthogonal with respect to
this uniform measure, the uniform measure is the natural \emph{prior}: it
expresses that we have no information, every direction on the sphere being
equally likely. Any other orthogonal state --- a definite Gegenbauer (and, in
the limit, Hermite) excitation --- encodes \emph{information}, namely the
result of integrating an unknown field's wavefunction over the hyperspherical
surface against the basis function attached to a point of that surface.

\subsection{The Gaussian limit of the hyperspherical weight}

Rescale the coordinate as $x\mapsto\sqrt{2/\alpha}\,x$ and let the dimension
$\alpha\to\infty$. The single-coordinate hyperspherical weight tends to the
Gaussian weight:
\begin{equation}
  \lim_{\alpha\to\infty}\Bigl(\sqrt{1-\tfrac{2x^2}{\alpha}}\Bigr)^{\alpha-1}
  =e^{-x^2}.
  \label{eq:gaussweight}
\end{equation}
Thus the uniform measure on an infinite-dimensional sphere is the Gaussian
measure --- the precise sense in which the Mehler ground state is the
``maximum-entropy / total-uncertainty'' prior, and in which the Gaussian source
$p_0=\mathcal{N}(0,\mathbf I)$ of a flow model is the Fock vacuum used in
Section~\ref{sec:framework}.

\subsection{Gegenbauer \texorpdfstring{$\to$}{to} Hermite and the normalization}

The orthogonal polynomials track the same limit. With the Gegenbauer
polynomials $C_n^{(\gamma)}$ defined by the generating function
$(1-2xw+w^2)^{-\gamma}=\sum_n C_n^{(\gamma)}(x)\,w^n$, L\'opez and
Temme~\cite{LopezTemme1999} give the Hermite limit (their Eq.~(1.4), with
$\gamma=\alpha/2$):
\begin{equation}
  \lim_{\alpha\to\infty}\Bigl(\tfrac{\alpha}{2}\Bigr)^{-n/2}
  C_n^{(\alpha/2)}\!\Bigl(\sqrt{\tfrac{2}{\alpha}}\,x\Bigr)
  =\frac{1}{n!}\,H_n(x),
  \label{eq:gegtoherm}
\end{equation}
where the Hermite polynomials are
$H_n(x)=n!\sum_{k=0}^{\lfloor n/2\rfloor}\frac{(-1)^k}{k!(n-2k)!}(2x)^{n-2k}$.
The normalizations converge consistently. For integers $n,\alpha\ge0$,
\begin{equation}
  N^{(\alpha)}_n
  =\Bigl(\tfrac{\alpha}{2}\Bigr)^{\frac12-n}
  \int_{-1}^{1}\bigl[C_n^{(\alpha/2)}(x)\bigr]^2
  \bigl(\sqrt{1-x^2}\bigr)^{\alpha-1}\,\mathrm{d}x
  =\Bigl(\tfrac{\alpha}{2}\Bigr)^{\frac12-n}
  \frac{\pi\,2^{1-\alpha}\,\Gamma(n+\alpha)}
       {n!\,\bigl(n+\tfrac{\alpha}{2}\bigr)\,[\Gamma(\tfrac{\alpha}{2})]^2},
  \label{eq:gegnorm}
\end{equation}
and, in the limit, this is exactly the (squared) Hermite normalization against
the Gaussian weight:
\begin{equation}
  \lim_{\alpha\to\infty}N^{(\alpha)}_n
  =\frac{\sqrt{\pi}\,2^n}{n!}
  =\int_{-\infty}^{\infty}\Bigl[\frac{H_n(x)}{n!}\Bigr]^2 e^{-x^2}\,\mathrm{d}x .
  \label{eq:hermnorm}
\end{equation}
(The last equality uses the standard orthogonality
$\int H_m H_n\,e^{-x^2}\mathrm{d}x=\sqrt{\pi}\,2^n n!\,\delta_{mn}$.)

\subsection{Reading the dictionary}

Equations~\eqref{eq:gaussweight}--\eqref{eq:hermnorm} establish a complete
dictionary, summarized in Table~\ref{tab:dictionary}. The infinite-dimensional
hypersphere with its uniform measure \emph{is} the Gaussian/Fock setting; the
hyperspherical harmonics \emph{are} the Hermite/number-state basis; the
uniform prior \emph{is} the vacuum $\ket{0}$; and the act of measuring a
component is the projection onto a definite excitation. This is exactly the
structure the \texttt{unfer} kernel manipulates natively in
\texttt{nested\textunderscore fock\textunderscore algebra}: the vacuum
projector $\ket{0}\bra{0}$ is the Gaussian source qsample $\Psi_0$, and creating
a boson in mode $j$ records one unit of information about data channel $j$.

\begin{table}[htbp]
\centering
\begin{tabular}{l|l}
\textbf{Finite hypersphere ($\alpha<\infty$)} & \textbf{Infinite-dim.\ limit (Fock)} \\
\hline
uniform surface measure $\mathrm{d}\sigma$ & Gaussian measure $e^{-x^2}\mathrm{d}x$ \\
Gegenbauer weight $(1-x^2)^{(\alpha-1)/2}$ & Gaussian weight $e^{-x^2}$
  \;[Eq.~\eqref{eq:gaussweight}] \\
Gegenbauer polynomials $C_n^{(\alpha/2)}$ & Hermite polynomials $H_n/n!$
  \;[Eq.~\eqref{eq:gegtoherm}] \\
hyperspherical harmonics & number states $\ket{n}$ \\
uniform ``no-information'' prior & Gaussian source qsample $\Psi_0=\ket{0}$ \\
a definite point / harmonic & an orthogonal excitation (information) \\
\end{tabular}
\caption{The hypersphere-to-Fock dictionary underlying the Gaussian source and
the vacuum prior.}
\label{tab:dictionary}
\end{table}

\subsection{Which vacuum? Mehler builds the outer Fock space's basis
directly}
\label{sec:innerouter}

Mehler's correspondence needs a space of wave-functions to act on, hence
needs the \emph{inner} Fock space (\texttt{nested\textunderscore
fock\textunderscore algebra}'s Level 1) to start from: the set of normalized
inner wave-functions is the infinite-dimensional hypersphere, and the
natural, rotation-invariant uniform measure on it expresses total ignorance
about which inner wave-function is present.
Equations~\eqref{eq:gaussweight}--\eqref{eq:hermnorm} show that, in the
infinite-dimensional limit, this uniform hyperspherical measure becomes the
Gaussian measure, and --- the point that matters here --- the Gegenbauer
polynomials $C_n^{(\alpha/2)}$ that are orthogonal with respect to it become
the Hermite polynomials $H_n$.

The $H_n$ are not merely a basis for functions on the limiting measure: by
the standard correspondence between the quantum harmonic oscillator and
Hermite functions, they \emph{are} the number-state basis of a bosonic Fock
space --- $H_0=1$, paired with the Gaussian weight, is the ground
state/vacuum $\ket{0}$, and each successive $H_n$ is the $n$-particle
excitation. Mehler's theorem therefore does not merely supply a measure to
be reinterpreted by some further isomorphism: \emph{it directly constructs
the vacuum and every other basis vector of a Fock space}, via the
Gegenbauer-to-Hermite limit. This constructed space is the \emph{outer} Fock
space: its one-particle ($H_1$-order) sector is exactly
$\ket{x_j}=\hat c_j^\dagger\ket{0}$, one per data point, matching
\texttt{nested\textunderscore fock\textunderscore algebra}'s representation
of an \texttt{OuterState} as a multiset of inner-space configurations
(\texttt{InnerBosonicState}/\texttt{InnerFermionicState}). Mehler applies to
the outer Fock space precisely because that is the space whose basis it
builds; the inner space only supplies the wave-functions --- the points of
the hypersphere --- that the construction needs to start from.

This is exactly why Table~\ref{tab:dictionary} may identify ``the vacuum
$\ket{0}$'' with the uniform prior over inner wave-functions: they are the
same object, $H_0$, produced by the one construction (Mehler's
Gegenbauer-to-Hermite limit) that builds the outer vacuum in the first
place --- not two vectors in two different spaces that happen to agree.

How, then, is this vacuum related to the data channels? Here the
parametrization of the next subsection does real work. A data packet
$\ket{x_j}$ is \emph{not} one of the Hermite/Gegenbauer basis elements: it
is a localized bump. In each of the $D$ coordinates that encode the point
$x^{(j)}$, its qsample is $\sqrt{1/w_i}$ on an arc of width $w_i$ of that
coordinate's circle (zero elsewhere), while in every remaining coordinate
it coincides with the uniform state. Expanding coordinate by coordinate in
the Gegenbauer/Hermite basis, the bump's component along the order-0
(vacuum) element of a localized coordinate is
$\int_{\mathrm{arc}}\sqrt{1/w_i}\,\sqrt{1/2\pi}\,\mathrm{d}\varphi
=\sqrt{w_i/2\pi}$, and each untouched coordinate contributes a factor $1$.
Hence
\begin{equation}
  \braket{0|x_j}=\varepsilon_j
  =\prod_{i=1}^{D}\sqrt{\frac{w_{j,i}}{2\pi}}\;>\;0 ,
  \label{eq:mehleroverlap}
\end{equation}
strictly positive \emph{precisely because} only finitely many coordinates
are localized: had infinitely many coordinates been disturbed, the infinite
product of factors $<1$ would vanish --- Kakutani's dichotomy for infinite
product measures. The uniform vacuum is therefore \emph{not} orthogonal to
any data channel. Distinct channels, by contrast, remain exactly
orthogonal, $\braket{x_i|x_j}=0$ for $i\neq j$: they localize the
\emph{same} $D$ coordinates on \emph{disjoint} arcs, so at least one
coordinate contributes the factor $\int\sqrt{p_i\,p_j}=0$ --- the
disjointness~\eqref{eq:disjoint} again.

For computation, \texttt{nested\textunderscore fock\textunderscore
algebra}'s \texttt{OuterState} basis is an \emph{orthonormal frame}: the
packet $\ket{x_j}$ is represented by the frame vector
$B_j^\dagger\ket{\mathrm{vac}}_F$ (the channels are orthonormal among
themselves, so they may be), and $\ket{\mathrm{vac}}_F$ is the unit vector
along the part of the uniform state lying \emph{outside} every packet ---
the Gram--Schmidt residual of $\ket{0}$ against the channels. In this frame
the Mehler vacuum is the \emph{dressed} superposition
\begin{equation}
  \ket{0}=c_0\ket{\mathrm{vac}}_F
  +\sum_{j=1}^M\varepsilon_j\,B_j^\dagger\ket{\mathrm{vac}}_F ,
  \qquad
  c_0=\sqrt{1-\textstyle\sum_j\varepsilon_j^2}\;,
  \label{eq:dressedvac}
\end{equation}
where $\sum_j\varepsilon_j^2\le1$ holds automatically: each
$\varepsilon_j^2$ is the uniform-measure mass of packet $j$'s support box,
and the boxes are disjoint. The code's rank-1
\texttt{Operator::ProjectVacuum} projects onto $\ket{\mathrm{vac}}_F$
alone, so identifying it with the Mehler prior $\ket{0}\bra{0}$ is exact
only to $\mathcal{O}(\varepsilon)$; the exact projector --- which
Section~\ref{sec:impl} shows is by itself the entire off-diagonal
generator --- is the \texttt{qfm\textunderscore hamiltonian\textunderscore
mehler\textunderscore projector} builtin.

\subsection{The data-channel wave-function on the hypersphere: finitely many
localized coordinates, the rest uniform}

Section~\ref{sec:framework} encodes each training point $x^{(j)}\in\Rset^D$
as a localized wave-packet $\Psi_j$ --- a single point of the inner Fock
space, hence a single point of the hypersphere of
Section~\ref{sec:mehler}. Concretely, in the finite-$k$ hypersphere of
angles $\varphi_0,\dots,\varphi_{k-1}$ ($D$ is fixed and finite, while
$k\to\infty$ in Mehler's limit), $\Psi_j$'s parametrization pins only the
first $D$ angular coordinates to a small, localized neighborhood of the
values encoding $x^{(j)}$'s $D$ real components; every remaining coordinate
$\varphi_D,\varphi_{D+1},\dots$ is left at the plain uniform measure on its
own circle, carrying no information about which data point this is ---
exactly as it does for the vacuum.

This is the hypersphere-side picture of ``one Hermite excitation, $D$ real
parameters, attached to point $x^{(j)}$,'' and it is what makes the encoding
a genuine \emph{single}, normalizable Fock-space excitation rather than an
infinite-dimensional one: a state can only differ from the vacuum in
finitely many modes, and localizing exactly $D$ of the (infinitely many)
hyperspherical coordinates is that finiteness condition made explicit. The
rest of the infinite-dimensional background stays at the uniform,
no-information state that Mehler identifies with the outer vacuum
(Section~\ref{sec:innerouter}) --- so $\ket{x_j}$ and $\ket{0}$ differ only
in those $D$ localized coordinates, and agree (uniform) everywhere else.
That agreement everywhere else is exactly what keeps the vacuum--channel
overlap~\eqref{eq:mehleroverlap} finite and nonzero: only the $D$ localized
coordinates contribute factors $\sqrt{w_i/2\pi}<1$, so
$\braket{0|x_j}=\varepsilon_j>0$ --- the data channels are visible to the
vacuum, and Section~\ref{sec:impl} turns precisely this non-orthogonality
into the off-diagonal transport.

\section{Krylov--Hashimoto dimensional reduction}
\label{sec:krylov}

The generator $\bar H$ of~\eqref{eq:Hbar} (and its diagonal
surrogate~\eqref{eq:Hdiag}) exactly encodes the flow, but its nominal dimension
scales with $M$. Where Layden et al.~\cite{Layden2025} discretize time with
product formulas (incurring $\|H_t^c\|$-dependent cost), \texttt{unfer}
compresses the generator with the SIRK machinery of
Section~\ref{sec:sirk}.

\subsection{The inversion-free shortcut}

The standard polynomial Krylov subspace is
\begin{equation}
  \Kry_m(\bar H,v_0)=\mathrm{span}\{v_0,\bar H v_0,\bar H^2 v_0,\dots,
  \bar H^{m-1}v_0\}.
\end{equation}
Naively, the rational construction~\eqref{eq:ratkrylov} requires solving a
linear system $(\gamma I-\bar H)y=v$ at each step. Our shortcut avoids it: by
pre-conditioning the seed as $v=(\gamma I-\bar H)^{m-1}v_0$, applying the
resolvent merely \emph{lowers the polynomial degree}, recovering the standard
Krylov subspace with no inversion. Concretely, \texttt{fock\textunderscore
sirk} realizes this by generating the sequence $w_k=(\bar H-z_kI)w_{k-1}$ with
uniform shifts $z_k=\gamma$; this is mathematically valid because $\bar H$ is
\emph{bounded} (the Mehler projector is rank one and the number operators are
diagonal). Each application of $\bar H^k$ costs only $\mathcal{O}(M)$ precisely
because the orthogonal Fock potential has \emph{no cross-terms} --- the same
decoupling that made the training~\eqref{eq:alpha} diagonal.

\subsection{Spectral low-pass filtering}

The Krylov projection acts as an exact spectral low-pass filter: it discards
the overfitted, localized ``spikes'' (high-frequency packet detail) and
retains the $m$ dominant macroscopic transport eigenvalues. The reduced
generator $\bar H_{\mathrm{reduced}}$ is a dense $m\times m$ matrix (e.g.\
$m=15$--$20$), and the SIRK bound~\eqref{eq:sirkbound} guarantees its error
decays like $e^{-hm}$. Generation is then a single instantaneous
$\mathcal{O}(m^2)$ matrix exponential,
\begin{equation}
  \Psi_{t=1}=e^{-i\bar H_{\mathrm{reduced}}}\,\Psi_0 .
\end{equation}
The orthonormalization of the (possibly rank-deficient) Krylov basis is handled
by \texttt{unfer}'s rank-revealing Gram whitening, which keeps the reduction
numerically stable even when several packets nearly coincide.

\section{Implementation in the \texttt{unfer} Kernel}
\label{sec:impl}

The wavefunction-flow module integrates into the existing \texttt{unfer}
monorepo with no new numerical primitives --- it reuses the CAS, the solver, and
the module runtime.

\subsection{Symbolic algebra (\texttt{nested\textunderscore fock\textunderscore algebra})}

The symbolic engine gains two operators. \texttt{Operator::ProjectVacuum}
is the rank-1 projector onto the \emph{frame} vacuum
$\ket{\mathrm{vac}}_F\bra{\mathrm{vac}}_F$ --- the
$\mathcal{O}(\varepsilon)$ stand-in for the exact Mehler source/prior
$H_0=\ket{0}\bra{0}$ of~\eqref{eq:dressedvac}; on application it keeps only
the strict-vacuum component and annihilates any state carrying a mode, and
it is self-adjoint and idempotent, so it slots into the existing adjoint
and Hermiticity machinery. \texttt{Operator::ProjectOnto($\psi$)} (rev 31)
is the exact rank-1 projector $\ket{\psi}\bra{\psi}$ onto an
\emph{arbitrary} normalized state --- in particular the dressed Mehler
vacuum of~\eqref{eq:dressedvac} --- applied through the rank-1 shortcut
$H\ket{s}=\braket{\psi|s}\ket{\psi}$ (one sparse inner product plus one
scaled copy), which is what makes the \emph{exact} off-diagonal generator
the practical one: cost $\mathcal{O}(M)$ per matvec, never the
$\mathcal{O}(M^2)$ frame expansion. The diagonal generator~\eqref{eq:Hdiag}
is instantiated directly by the builtin
\texttt{qfm\textunderscore hamiltonian(alphas)}, which assembles
$\ket{0}\bra{0}+\sum_j\bar\alpha_j\,\hat c_j^\dagger\hat c_j$ as one projector
term plus one number operator per data point. Building the Hamiltonian
directly (rather than through \texttt{Expression::expand()}) sidesteps the
combinatorial explosion of the CAS, so $M=10^6$ data points yield $10^6$ terms
with no blow-up --- the grid-free analogue of the paper avoiding an explicit
$N^d$ amplitude vector. Both \texttt{qfm\textunderscore hamiltonian} and the
exact projector builtin (below) identify each data point purely by array
index, which already guarantees $\braket{x_i|x_j}=\delta_{ij}$ but carries
none of the point's own $D$ real coordinates in the Fock state itself, only
in the scalar weights. The companions
\texttt{qfm\textunderscore hamiltonian\textunderscore localized}/
\texttt{qfm\textunderscore hamiltonian\textunderscore mehler\textunderscore
projector\textunderscore localized} instead realize
Section~\ref{sec:mehler}'s data-channel picture literally: each
\texttt{point\textunderscore to\textunderscore inner\textunderscore state}
occupies exactly $D$ inner modes, one per real coordinate (a fixed-point
quantization, zigzag-encoded to preserve sign), leaving every other mode at
zero occupation --- the direct computational realization of ``$D$ of the
hyperspherical coordinates localized, the rest uniform,'' reachable through
the kernel's protocol as the \texttt{qfm\textunderscore mehler\textunderscore
localized}/\texttt{qfm\textunderscore mehler\textunderscore
projector\textunderscore localized} builtins.

\subsection{GPU-accelerated solver (\texttt{fock\textunderscore sirk})}

The solver natively generates the Krylov sequence $w_k=(\bar H-z_kI)w_{k-1}$.
Supplying a uniform shift array realizes the inversion-free reduction of
Section~\ref{sec:krylov}; the $m\times m$ Gram matrix $\braket{w_j|w_k}$ is
assembled on the GPU via \texttt{candle-core} CUDA tensors. Because
$\hat V=\sum_j\bar\alpha_j\hat n_j$ has no cross-terms, basis generation is
embarrassingly parallel and fast, and the rank-revealing Gram whitening keeps
the highly compressed flow representation stable. This is \texttt{unfer}'s
classical counterpart to the paper's quantum Hamiltonian simulation: exact
operator evolution, but on a CPU/GPU Fock substrate rather than qubits.

\subsection{Module execution runtime}

The whole wavefunction-flow loop is exposed through the
\texttt{unfer\textunderscore ffi} C ABI and driven by a new Austral module
(\texttt{qfm\textunderscore module}). Through the manifest-authorized
\texttt{uk\textunderscore evolve} and
\texttt{uk\textunderscore event\textunderscore probability} symbols, the module
triggers the analytical parameter assembly, the Krylov reduction, and the
single-step $\mathcal{O}(m^2)$ inference natively in-process, with the kernel
handle held in a linear type that the Austral compiler forces it to free. A
second module, \texttt{data\textunderscore source}, exercises the
\texttt{data\textunderscore source} archetype by linking
\texttt{unfer\textunderscore ffi} as a path dependency and ingesting
external observations through the real C ABI (\texttt{uk\textunderscore observe},
\texttt{uk\textunderscore get\textunderscore result}); it runs as
\texttt{step~8} of \texttt{demo\textunderscore module/run\textunderscore demo.sh}.

\subsection{Scope: the diagonal surrogate and the exact off-diagonal projector}

The shipped builtin \texttt{qfm\textunderscore hamiltonian(alphas)} implements
the diagonal generator~\eqref{eq:Hdiag}. Because $\bar H_{\text{impl}}$ is
diagonal, the vacuum and the data channels are eigenstates, so
$e^{-i\bar H_{\text{impl}}t}$ contributes only phases and the Born populations
are stationary --- the module's integration tests assert exactly this
conservation, not a spurious ``spread.'' This is the spectrally-stable
skeleton of~\eqref{eq:Hbar}: it shares the prior, the $\mathcal{O}(M)$
construction, and the SIRK compressibility, while the off-diagonal mixing
$\hat h_j$ --- the part that actually performs the wavefunction flow --- is
taken up next.

\paragraph{The exact off-diagonal generator is just the vacuum projector.}
Section~\ref{sec:innerouter} established the key structural fact: the
Mehler vacuum is \emph{not} orthogonal to the data channels ---
$\braket{0|x_j}=\varepsilon_j>0$~\eqref{eq:mehleroverlap}, because each
channel localizes only finitely many hyperspherical coordinates. The exact
off-diagonal generator therefore needs no explicit coupling terms at all:
it is \emph{just the projector onto the vacuum},
\begin{equation}
  \Hgen_{\mathrm{off}}=\ket{0}\bra{0},\qquad
  \bra{x_i}\Hgen_{\mathrm{off}}\ket{x_j}
  =\varepsilon_i\varepsilon_j ,
  \label{eq:Hmehler}
\end{equation}
off-diagonal in the channel basis purely by non-orthogonality. It is
manifestly Hermitian (every projector is), so it drops into the SIRK +
Born-rule substrate with no symmetrization step. For analytical reference,
substituting the dressed representation~\eqref{eq:dressedvac} expands it
over the orthonormal computational frame as
\begin{equation}
  \ket{0}\bra{0}
  =c_0^2\,P_0
  +\sum_{j=1}^M c_0\varepsilon_j\bigl(B_j^\dagger P_0+P_0B_j\bigr)
  +\sum_{i,j=1}^M \varepsilon_i\varepsilon_j\,B_i^\dagger P_0 B_j ,
  \label{eq:Hmehlerexp}
\end{equation}
where $P_0=\ket{\mathrm{vac}}_F\bra{\mathrm{vac}}_F$ is the \emph{frame}-vacuum
projector (\texttt{Operator::ProjectVacuum}) and
$B_j^\dagger=\hat C^\dagger_{\ket{x_j}}$ creates the single-excitation
outer-universe state --- but the expansion is \emph{never materialized}:
the implementation stores the dressed vacuum once and applies the rank-1
shortcut
\begin{equation}
  \Hgen_{\mathrm{off}}\ket{s}=\braket{0|s}\,\ket{0}
  \qquad(\texttt{Operator::ProjectOnto}),
  \label{eq:rank1shortcut}
\end{equation}
one sparse inner product plus one scaled copy, $\mathcal{O}(M)$ per matvec
--- so no $\mathcal{O}(\varepsilon)$ truncation and no
$\mathcal{O}(M^2)$ cross-term list is ever needed. (Earlier revisions
shipped a truncated companion builtin, $\ket{0}\bra{0}+\sum_j\bar\alpha_j
(B_j^\dagger P_0+P_0B_j)$ with $c_0\to1$ and the cross terms dropped; rev
31 removed it --- the exact projector is the only off-diagonal generator.)

Because the exact generator is a rank-1 projector ($H^2=H$, eigenvalues
$1$ on the dressed $\ket{0}$ and $0$ on its orthogonal complement), its
evolution is closed-form,
\begin{equation}
  e^{-it\,\Hgen_{\mathrm{off}}}
  =I+\bigl(e^{-it}-1\bigr)\ket{0}\bra{0}\,.
\end{equation}
Starting from the frame vacuum $\ket{\mathrm{vac}}_F$ (overlap $c_0$ with
$\ket{0}$), \emph{every} channel is pumped coherently and simultaneously at
the same unit frequency, $P_{x_j}(t)=4\sin^2(t/2)\,c_0^2\varepsilon_j^2$,
with exact return at $t=2\pi$; the dressed vacuum itself is stationary (it
is the eigenvector). This coherent, reversible oscillation --- rather than
the diffusive Fokker--Planck transport of the paper's anti-Hermitian
$\Hgen^c$~\cite{Layden2025} --- is the honest engineering trade required
by the unfer architecture: SIRK + the Born rule assume a Hermitian
generator (AGENTS.md \S4). The generator is implemented by
\texttt{qfm\textunderscore hamiltonian\textunderscore mehler\textunderscore
projector(epsilons)} in \texttt{nested\textunderscore fock\textunderscore
algebra} --- with $\varepsilon_j$ obtainable from the arc widths via
\texttt{mehler\textunderscore channel\textunderscore overlap(widths)}
$=\prod_i\sqrt{w_i/2\pi}$ --- and dispatched as the
\texttt{qfm\textunderscore mehler\textunderscore projector} builtin in
\texttt{prob\textunderscore kernel::build} (which rejects
$\sum_j\varepsilon_j^2>1$ with a diagnostic); the localized-encoding
variant is \texttt{qfm\textunderscore mehler\textunderscore
projector\textunderscore localized(points, epsilons, scale)}. Unit tests
verify idempotence $H^2=H$ on arbitrary states and every matrix element of
\eqref{eq:Hmehler}; integration tests pin the SIRK-evolved populations to
the closed form above and the stationarity of the dressed vacuum.

\paragraph{Beyond off-diagonal: Tomographic Subspace Recovery (Workstream F).}
The diagonal and exact-projector generators both encode each of the
$M$ data points as a single-boson mode of the Fock space, so online
generation is just $e^{-iH_m t}$ on the Krylov-reduced $c_0$ --- there is
no $M$-independent decoding of a raw high-resolution image. Section~\ref{sec:tomographic}
introduces the \emph{Tomographic Subspace Recovery} (TSR) pipeline that closes
this gap: a two-level hashing scheme plus pre-projected observables in the
$m^2$ operator basis that execute the full four-phase generate
(encode $\to$ evolve $\to$ tomographic reconstruct $\to$ lossless decode) with
zero $\mathcal{O}(M)$ dependence at inference. TSR is the natural v2 of QFM
that the unfer architecture now supports.

\section{Tomographic Subspace Recovery: a non-neural,
\texorpdfstring{$M$}{M}-free online pipeline}
\label{sec:tomographic}

The diagonal and exact-projector generators of
Sections~\ref{sec:framework}--\ref{sec:impl} each encode the $M$ data points
as single-boson modes of the Fock space. Online generation is then
just $e^{-iH_m t}$ on the Krylov-reduced coefficient vector $c_0\in\Cset^m$;
there is no mechanism to decode a raw high-resolution image $x\in\Rset^d$
back from the Krylov state. This section closes that gap by adapting the
algorithm specification \emph{Non-Neural Quantum Flow Matching with Tomographic
Subspace Recovery} (henceforth TSR) into the unfer architecture. TSR
decouples \emph{semantics} (hashing and coordinate projections) from
\emph{reasoning} (unitary Krylov evolution) so that every online op is
$\mathcal{O}(d\cdot m^2)+\mathcal{O}(K_2\log k)+\mathcal{O}(K_2\cdot m^2)$
with \emph{no} $\mathcal{O}(M)$ dependence.

\subsection{System dimensions}

The pipeline is defined over the following five dimensions:
\begin{itemize}
\item $d$: raw configuration space dimension (e.g.\ $d=10^6$ pixels);
\item $M$: total number of training data points (offline only);
\item $k$: Level~1 spatial feature hash dimension (e.g.\ $k=512$, $k\ll d$);
\item $K_2$: Level~2 sketched Hilbert space dimension (e.g.\ $K_2=8192$, $K_2>k$);
\item $m$: reduced Krylov subspace dimension (e.g.\ $m=20$).
\end{itemize}
The two-level hashing is the key compression step: $S_1$ reduces the raw
configuration to a $k$-dim feature vector, and $S_2$ maps each feature to a
single-excitation Fock state in a $K_2$-dim space. After Krylov reduction the
online state lives in $\Cset^m$, and all raw-coordinate observables are
pre-projected into the $m^2$ operator basis so that the full four-phase
generate costs are bounded by the four combinations above.

\subsection{Offline compilation
(\texorpdfstring{$\mathcal{O}(M)$}{O(M)})}

The offline phase builds the static artifacts that the online phase consumes.

\paragraph{Level~1 hash ($S_1\in\Rset^{k\times d}$).} A sparse, deterministic
Count-Sketch matrix maps each raw pixel coordinate $c\in\{1,\dots,d\}$ to a
hash bucket $h(c)\in\{1,\dots,k\}$ and a random sign $s(c)\in\{-1,+1\}$. The
full $d\times k$ matrix is never materialized --- only the per-coordinate
$(h(c),s(c))$ pairs are stored ($2d$ bytes). For any image $x\in\Rset^d$ the
spatial hash is
\begin{equation}
  \tilde x=S_1(x)\in\Rset^k,\qquad \tilde x_h=\sum_{c:h(c)=h}s(c)\,x_c .
\end{equation}
Complexity: $\mathcal{O}(\mathrm{nnz}(x))$ for a sparse $x$, $\mathcal{O}(d)$
dense. In the unfer implementation this is
\texttt{CountSketch::new(k,d,seed)} in \texttt{qfm/src/sketch.rs} with a
splitmix64 PRNG over the seed.

\paragraph{Level~2 hash ($S_2:\mathcal{P}(\Rset^k)\to\Cset^{K_2}$).} The
Level~2 hash maps a $k$-dim feature distribution to a $K_2$-dim sketched
Fock state. For a single training image $x_j$ the distribution is a delta
function $\delta_{\tilde x_j}$ at the feature point; the corresponding Fock
state is a \emph{single excitation} at the mode assigned to that feature:
\begin{equation}
  \ket{\Psi_j}=S_2(\delta_{\tilde x_j})\in\Cset^{K_2},\qquad
  \ket{\Psi_j}=B^\dagger_{m_j}\ket{0}\quad\text{(mode }m_j\text{)} .
\end{equation}
In unfer this is \texttt{FeatureToMode::register(feature\_key)} which assigns
a fresh mode to each new training feature and stores the
$(\text{feature\_key}\to\text{mode})$ map in an \texttt{FxHashMap}.

\paragraph{Analytical potential and exact flow generator.} The
conservative potential $V_t(x)=\sum_j\alpha_j(t)\Psi_j(x)$ with the
disjoint-support $\Psi_j$ of Section~\ref{sec:framework} gives the
decoupled, linear-scaling optimal coefficients of~\eqref{eq:alpha}. The
time-averaged coefficients $\bar\alpha_j=\int_0^1\alpha_j(t)\,\mathrm{d}t$
define a static, time-independent flow potential. In the unfer
implementation the optimal coefficients reduce to the closed form
\begin{equation}
  \bar\alpha_j=\frac{\|x^{(j)}\|^2}{M},
\end{equation}
the squared norm of each data point divided by the dataset size (a positive
weight proportional to the displacement energy of the $j$-th channel).
The static flow generator is the exact Mehler
projector~\eqref{eq:Hmehler}, with the flow-matching weights setting the
channel components of the dressed vacuum: normalizing the dressed vector
$\ket{v}=\ket{\mathrm{vac}}_F+\sum_{j=1}^{K_2}\bar\alpha_j\ket{x_j}$ gives
\begin{equation}
  \bar H_{\mathrm{tomo}}=\ket{\tilde 0}\bra{\tilde 0},\qquad
  \ket{\tilde 0}=\frac{\ket{v}}{\|v\|}
  =c_0\ket{\mathrm{vac}}_F+\sum_{j=1}^{K_2}\varepsilon_j\ket{x_j},\qquad
  \varepsilon_j=\frac{\bar\alpha_j}{\sqrt{1+\sum_k\bar\alpha_k^2}},
  \label{eq:Htomo}
\end{equation}
with $c_0=1/\sqrt{1+\sum_k\bar\alpha_k^2}$ --- so
$\varepsilon_j\to\bar\alpha_j$ for small weights, and the generator is
exactly idempotent for any weights. Only channels $j<K_2$ (the sketched
modes) enter the dressed vacuum. The vacuum couples to all $K_2$ channels
simultaneously through the single projector (this is not $K_2$ independent
two-level oscillators), transporting amplitude coherently with the
closed-form populations of Section~\ref{sec:impl}. This is implemented by
\texttt{qfm/src/potential.rs::build\_flow\_hamiltonian}, which delegates
to \texttt{qfm\textunderscore hamiltonian\textunderscore mehler\textunderscore
projector} and therefore to the rank-1
\texttt{Operator::ProjectOnto} shortcut~\eqref{eq:rank1shortcut}.

\paragraph{Krylov reduction.} Running the SIRK pipeline of
Section~\ref{sec:krylov} on $\bar H_{\mathrm{tomo}}$ with a vacuum seed
produces the $K_2\times m$ basis matrix $W$ and the $m\times m$ reduced
generator $H_m$. In the unfer implementation, \texttt{QfmPipeline::compile}
calls \texttt{fock\_sirk::solve\_forward\_sirk} on the vacuum $+$
single-excitation seed ($K_2+1$ basis states, each with amplitude
$1/\sqrt{K_2+1}$) with $m$ uniform shifts on the negative-imaginary axis
($\zeta_k=-i\,k/m$ for $k=1,\dots,m$; the shifts are normalized by $m$
since rev 24 to bound the per-step growth of the Krylov sequence --- see
the rev 24 subsection of Section~\ref{sec:bayes} for the overflow this
fixes). The SIRK solve returns a Hermitian reduced Hamiltonian
$H_m=\mathrm{sirk.h\_proj}$ (Hermiticity is guaranteed by the Gram-matrix
whitening step in the forward SIRK pipeline of Section~\ref{sec:krylov}) plus a
whitened Krylov basis $W=\mathrm{sirk.w\_whiten}$. Since rev 17
(\texttt{P6 G}) the spatial mode basis used for encode/decode is this
genuine SIRK-generated $\mathrm{w\_whiten}$, restricted to the $K_2$
single-excitation Fock rows and per-row renormalized so the Born-rule
likelihood stays well-defined. (The rev 14 release had interim-materialized
$W$ as the deterministic identity sub-block of the first $m$ standard basis
vectors --- a valid rank-$m$ stand-in exploiting the spec's \emph{``$K_2$-dim
single-excitation subspace is small enough for direct construction''}
insight --- but this left a lossy component in the round-trip; rev 17
replaced it.) The reduced Hamiltonian
$H_m$ is genuinely non-diagonal (the off-diagonal coupling comes
from the static flow generator itself, $\bar H_{\mathrm{tomo}}$ of
\eqref{eq:Htomo}), so Phase~2 below uses \texttt{nalgebra}'s
Pad\'e-based matrix exponential on $(-i H_m t)$ to preserve unitarity
exactly (AGENTS.md~\S4). The important point is that $W$ and $H_m$ are
\emph{compiled once} and then retained; the raw $M\times d$ training
dataset and the full $K_2$-dim Fock state vectors are purged from memory
at the end of compilation (only the $M$ sketched $k$-dim training features
are kept, for the $S_2$ nearest-feature fallback of the online phase).

\paragraph{Pre-projected observables.} The online phase never touches the
raw Fock space or the raw coordinates. Instead, all observables are
pre-projected into the $m^2$ operator basis
$\{E_{r,s}=\ket{e_r}\bra{e_s}:r,s\in\{1,\dots,m\}\}$ (a complete,
orthonormal basis of Krylov operators) and stored as dense static matrices.
Four matrices are built.

The \emph{probability weight matrix} $\mathbf{W}_{\mathrm{prob}}\in\Rset^{K_2\times m^2}$
projects the density matrix back to the sketched probability space. For each
basis projector $P_a=\ket{a}\bra{a}$ ($a=0,\dots,K_2-1$) and each
$E_{r,s}$:
\begin{equation}
  (\mathbf{W}_{\mathrm{prob}})_{a,(r,s)}
  =\mathrm{Tr}\bigl(E_{r,s}^\dagger\,W^\dagger P_a W\bigr)
  =(W^\dagger P_a W)_{s,r}
  =\overline{W_{a,r}}\,W_{a,s}.
\end{equation}
The last equality uses that $P_a$ is a one-hot projector: $W^\dagger P_a W$ is
the outer product of the $a$-th row of $W$ with itself. Implemented by
\texttt{qfm/src/observables.rs::probability\_weight\_matrix}; complexity
$\mathcal{O}(K_2\cdot m^2)$ (one outer product per projector).

The \emph{Krylov image basis} $\mathbf\Phi\in\Rset^{d\times m^2}$ projects
the raw coordinate (pixel) operators $X_c=\ket{c}\bra{c}$ ($c=0,\dots,d-1$)
into the Krylov subspace:
\begin{equation}
  \mathbf\Phi_{c,(r,s)}=\mathrm{Tr}\bigl(E_{r,s}^\dagger\,W^\dagger X_c W\bigr)
  =\overline{W_{c,r}}\,W_{c,s}.
\end{equation}
Implemented by \texttt{krylov\_image\_basis}; complexity $\mathcal{O}(d\cdot m^2)$.

The \emph{compressive subspace solver} $\tilde{\mathbf\Phi}^+\in\Rset^{m^2\times k}$
applies the Level~1 hash to $\mathbf\Phi$ and pre-computes its
Moore--Penrose pseudo-inverse:
\begin{equation}
  \tilde{\mathbf\Phi}=S_1\mathbf\Phi\in\Rset^{k\times m^2},\qquad
  \tilde{\mathbf\Phi}^+=(\tilde{\mathbf\Phi}^\top\tilde{\mathbf\Phi})^{-1}
  \tilde{\mathbf\Phi}^\top\in\Rset^{m^2\times k}.
\end{equation}
Implemented by \texttt{CountSketch::apply\_to\_columns} (the $S_1\mathbf\Phi$
step, $\mathcal{O}(k\cdot m^2)$) and \texttt{compressive\_solver} (SVD-based
pseudo-inverse, $\mathcal{O}(k\cdot m^4)$ dominated by the SVD).

The offline compilation is complete. The \texttt{QfmPipeline} struct retains
the $K_2\times m$ basis $W$, the reduced generator $H_m$, the three
pre-projected observable matrices, and the $M$ sketched $k$-dim training
features (for the $S_2$ nearest-feature fallback); the raw $M\times d$
training dataset is discarded, and every online op touches only these
compiled artifacts --- never the raw training data.

\subsection{Online inference: the four-phase generate}

The online phase has no $\mathcal{O}(M)$ dependencies. Given a raw query
$q\in\Rset^d$, the four phases below chain
\begin{equation}
  q\,\xrightarrow{S_1}\,\tilde q\,\xrightarrow{S_2}\,\ket{\Psi_{\mathrm{in}}}\,\xrightarrow{W^\dagger}\,c_0\,\xrightarrow{e^{-iH_m}}\,c_1\,\xrightarrow{\mathrm{tomo+decode}}\,x_{\mathrm{out}},
\end{equation}
and each phase is bounded by the complexity table at the end of this section.
In the unfer implementation these phases are the
\texttt{QfmPipeline::encode}/\texttt{evolve}/\texttt{decode} methods, with
\texttt{generate} chaining all three.

\paragraph{Phase 1: Input encoding.} The sparse, single-excitation Fock state
$\ket{\Psi_{\mathrm{in}}}=S_2(S_1(q))\in\Cset^{K_2}$ is projected into the
$m$-dim starting vector $c_0=W^\dagger\ket{\Psi_{\mathrm{in}}}$. Since the
state is single-excitation, only the $K_2$ active entries of the dense
$K_2$-dim vector contribute: complexity $\mathcal{O}(K_2\cdot m)$ reduced to
$\mathcal{O}(k\cdot m)$ by the Level~1 hash. In unfer this is
\texttt{QfmPipeline::encode} returning a \texttt{DVector<Complex64>} of
length $m$.

\paragraph{Phase 2: Unitary time-evolution.} The state evolves to the
generation boundary $t=1$ by $c_1=e^{-iH_m t}\,c_0$. Since rev 14 the
reduced Hamiltonian $H_m$ is genuinely non-diagonal (it is the SIRK
projection of the off-diagonally coupled $\bar H_{\mathrm{tomo}}$ of
\eqref{eq:Htomo}), so this is a full Pad\'e-based matrix exponential via
\texttt{nalgebra::DMatrix::exp()} applied to the skew-Hermitian matrix
$-i H_m t$. The Pad\'e approximant is provably unitary and stable for
Hermitian $H_m$, exactly matching the AGENTS.md~\S4 mandate for
inverse-free SIRK. The time parameter $t$ is taken explicitly by
\texttt{QfmPipeline::evolve(c\_0, t)} (with the default $t=1.0$ exposed
via the wrapper \texttt{generate(query)}). Complexity:
$\mathcal{O}(m^3)$ for the Pad\'e exponentiation; in practice this is
$<$30\,\textmu s on commodity hardware for $m\leq 32$, and the rev 14
test \texttt{pipeline\_evolve\_unitarity\_preserves\_norm} verifies
$|||c_0||-||c_1||| < 10^{-6}$ for arbitrary $t$. Implemented by
\texttt{QfmPipeline::evolve}.

\paragraph{Phase 3: Tomographic distribution reconstruction.} The density
matrix of $c_1$ is computed in the $m^2$ operator basis: the flat outer
product is $m^2$ complex multiplications, but only its real part is kept,
$\rho_{\mathrm{flat}}=\mathrm{Re}\,\mathrm{vec}(c_1c_1^\dagger)\in\Rset^{m^2}$
--- lossless here because the SIRK-whitened basis $W$ is real by
construction ($\bar H_{\mathrm{tomo}}$ is real-symmetric in this basis), so
$W_{\mathrm{prob}}$'s own projection weights are real too and the discarded
imaginary parts of $\rho_{\mathrm{flat}}$ would only have multiplied zero.
The sketched probability density is then
$\tilde p=\mathbf{W}_{\mathrm{prob}}\,\rho_{\mathrm{flat}}\in\Rset^{K_2}$:
a single dense matrix-vector product, complexity $\mathcal{O}(K_2\cdot m^2)$.
In unfer this is the first half of \texttt{QfmPipeline::decode}
(\texttt{qfm/src/observables.rs::probability\_weight\_matrix} builds
$W_{\mathrm{prob}}$ under the same real-$W$ assumption).

\paragraph{Phase 4: Lossless peak recovery (decoding).} The Count-Sketch
Heavy Hitters algorithm locates the top-$k$ indices of the probability sketch
$\tilde p$ (\texttt{HeavyHitters} in \texttt{qfm/src/heavy\_hitters.rs}, a
Misra--Gries / Count-Min tracker, complexity $\mathcal{O}(K_2\log k)$). The
peak mode's training feature is looked up to recover the $k$-dim peak hash
$\tilde x_{\mathrm{peak}}$. The subspace coefficients
$\gamma=\tilde{\mathbf\Phi}^+\,\tilde x_{\mathrm{peak}}\in\Rset^{m^2}$ are
recovered by a single small matrix-vector multiplication
($\mathcal{O}(m^2\cdot k)$). The full-resolution image is rendered as
$x_{\mathrm{out}}=\mathbf\Phi\,\gamma\in\Rset^d$ ($\mathcal{O}(d\cdot m^2)$).
In unfer this is the second half of \texttt{QfmPipeline::decode}.

\subsection{Complexity and resource scaling}

The total online inference cost is
\begin{equation}
  \mathrm{Total\;FLOPs}=\mathcal{O}(d\cdot m^2)+\mathcal{O}(K_2\log k)
  +\mathcal{O}(K_2\cdot m^2),
\end{equation}
with no $M$ dependence. The load-bearing complexity allocations are:

\begin{itemize}
\item \textbf{Image render} ($\mathbf\Phi\,\gamma$): $\mathcal{O}(d\cdot m^2)$ ---
  the dominant cost for high-resolution $d$. The output buffer is the
  bottleneck, not the compute.
\item \textbf{Sketched probability} ($\mathbf{W}_{\mathrm{prob}}\,\rho$):
  $\mathcal{O}(K_2\cdot m^2)$ --- a single BLAS \texttt{dgemv} call, sub-ms
  on commodity hardware.
\item \textbf{Heavy hitters} ($\mathrm{HeavyHitters}(\tilde p)$):
  $\mathcal{O}(K_2\log k)$ --- a binary-tree search over the top-$k$ tracker.
\item \textbf{Subspace recovery} ($\tilde{\mathbf\Phi}^+\,\tilde x$):
  $\mathcal{O}(m^2\cdot k)$ --- a small matrix-vector product that lives in
  L1 cache.
\item \textbf{Krylov projection} ($W^\dagger\ket{\Psi_{\mathrm{in}}}$):
  $\mathcal{O}(k\cdot m)$ --- only the $k$ active columns of $W^\dagger$ are
  touched (the single-excitation sparsity of $\ket{\Psi_{\mathrm{in}}}$).
\item \textbf{Unitary flow} ($e^{-iH_m}\,c_0$): $\mathcal{O}(m^2)$ ---
  instantaneous in L1 cache.
\item \textbf{Tomography} ($\mathrm{vec}(c_1c_1^\dagger)$): exactly $m^2$
  complex multiplications --- e.g.\ $400$ for $m=20$.
\end{itemize}

\subsection{Integration into unfer}

TSR is exposed through the existing unfer probability-kernel surface via a
new \texttt{HamiltonianSpec::QfmTomography} variant
(\texttt{unfer\_protocol/src/types.rs}) carrying the
\texttt{QfmTomographySpec} compilation spec. The
\texttt{prob\_kernel::Session} detects this variant in
\texttt{Session::new}, compiles a \texttt{QfmPipeline} from the spec, and
stores it alongside the placeholder Hamiltonian. The \texttt{evolve\_with\_query}
method (and the FFI \texttt{uk\_evolve} with an optional
\texttt{"query": [f64; d]} field) dispatches to the pipeline's 4-phase
\texttt{generate}; the result populates the new
\texttt{EvolveReport::qfm\_output} field, which is serialized as the
generated raw image. The error path (\texttt{KernelError::Qfm(..)}) carries
the QFM-specific UK-coded diagnostics with \texttt{RepairHint}s (e.g.\
dimension-mismatch: replace \texttt{evolve.query} with the right
dimensionality; degenerate basis: increase \texttt{hamiltonian.spec.krylov\_dim}).

\subsection{Honest scope limits}

\begin{itemize}
\item \textbf{Heavy hitters returns a single peak} (top-1). Multi-modal
generation (top-$k>1$) is a future extension.
\item \textbf{The $S_2$ nearest-feature fallback} for unseen queries is a
heuristic (L1 distance over the $k$-dim sketch). The spec's delta-function
mapping assumes every query hits a training point; the fallback is the
honest engineering compromise.
\item \textbf{Hermiticity deviation.} The spec writes the flow generator as
the anti-Hermitian $\bar H=\ket{0}\bra{0}-\frac{i}{2}\sum\bar\alpha_j\hat h_j$,
producing irreversible Fokker--Planck transport. unfer's SIRK + Born rule
require a Hermitian generator (AGENTS.md \S4), so the implementation uses
the exact Mehler projector
$\bar H_{\mathrm{tomo}}=\ket{\tilde 0}\bra{\tilde 0}$~\eqref{eq:Htomo}
--- coherent evolution, not diffusive transport.
This is the right call architecturally: the coherent form lets
$e^{-iH_m t}$ be a true unitary flow, exactly the AGENTS.md~\S4 mandate,
and the Pad\'e-based matrix exponential is provably unitary and
Hermiticity-preserving. The rev 14 test
\texttt{pipeline\_evolve\_unitarity\_preserves\_norm} verifies
$|||c_0||-||c_1||| < 10^{-6}$ for arbitrary $t$.
\item \textbf{Raw-coordinate basis.} The spec's $\mathbf\Phi$ assumes the
raw domain is a direct sum of one-hot pixel projectors $X_c=\ket{c}\bra{c}$.
A more general continuous-position model would require $X_c$ to be derived
from a discretized position basis --- the spec leaves this implicit.
\item \textbf{Spatial mode basis $W$ (resolved in rev 17, P6 G).} The
rev 14/15 releases materialized the encode/decode basis $W$ as the
deterministic identity sub-block of the first $m$ standard basis vectors
in the $K_2$-dim Fock space, which left a small lossy component in the
decompression round-trip at high $d$. Since rev 17, $W$ is the genuine
SIRK-generated $\mathrm{w\_whiten}\in\Cset^{K_2\times m}$ (restricted to
the $K_2$ single-excitation rows, per-row renormalized), at the cost of
$K_2\cdot m$ storage; the constraint $\texttt{krylov\_dim}\ge K_2$ that
makes the row restriction well-defined is enforced at runtime since
rev 18 (P7.4), and can be relaxed via the rev 26 rank-truncation path
(\texttt{max\_rank}, see the rev 25--28 subsection of
Section~\ref{sec:bayes}).
\end{itemize}

\subsection{Rev 15 outcomes: end-to-end module demo + benchmarks}

Two final v2 frontier items were resolved in rev 15 to give the TSR
pipeline the same module-level integration coverage as the rest of the
unfer surface.

\paragraph{\texttt{qfm\_tomo\_module/} Austral demo.} A new module
directory \texttt{unfer/qfm\_tomo\_module/} mirrors the existing
\texttt{qfm\_module/} pattern (\texttt{module.toml} +
\texttt{QfmTomoModule.aui}/\texttt{QfmTomoModule.aum} +
\texttt{build.sh} + \texttt{run\_demo.sh}). It JIT-creates a
\texttt{qfm\_tomography} model from the 4-point tetrahedron training set
used in the FFI integration tests, runs the 4-phase generate with a
\texttt{query}, drains the \texttt{EvolveReport} via
\texttt{uk\_get\_result}, and frees via the linear \texttt{Model}. The
UK-4001 authorization gate is exercised by the run script: granting
\texttt{uk\_get\_result} in the manifest produces \texttt{ALLOW};
revoking it produces \texttt{UK-4001 CallDenied}. The module is wired
into \texttt{.github/workflows/ci.yml} as the \texttt{qfm-tomo-e2e}
job, mirroring the existing \texttt{qfm-e2e} job.

\paragraph{Criterion benchmarks (\texttt{qfm/benches/pipeline.rs}).}
Three benchmark groups confirm the architecture's central scaling
claims. \texttt{compile\_vs\_M} (training-set size $M=10/100/1000$)
measures the offline $\mathcal{O}(M)$ cost; \texttt{generate\_vs\_d}
(raw dimension $d=64/256/1024$) measures the online $\mathcal{O}(d)$
image-render cost; \texttt{sketch\_apply\_vs\_d} (raw dimension
$d=64/256/1024/4096$) measures the Level~1 Count-Sketch's $\mathcal{O}(d)$
cost. Sample measurements on the local rev 15 run: compile($M=1000$)
$\approx 3.6$\,s, generate($d=1024$) $\approx 20$\,\textmu s,
sketch\_apply($d=4096$) $\approx 9$\,\textmu s --- all consistent with
the asymptotic complexity table above.

\section{Quantum Bayesian Updating on the TSR-evolved Prior}
\label{sec:bayes}

The TSR pipeline of \S\ref{sec:tomographic} evolves the uninformative
Mehler prior $\mu_0$ on the unit sphere of $\Cset^m$ through a unitary
flow $U_m=e^{-iH_m t}$ to the generation boundary $t=1$. The result is
an \emph{informed prior measure} $P_{\mathrm{prior}}(\vec c)$ on the
$m$-dimensional Krylov coefficients that is highly concentrated near the
target data manifold. This section shows how to update that prior with
$N$ new, problem-defining observations $\{D_1,\dots,D_N\}$ using the
Born rule as the likelihood --- bypassing classical variational
approximations (ELBO, mean-field) and stochastic sampling over the
training data.

The architecture is the same two-level hashing + Krylov projection
machinery of \S\ref{sec:tomographic}; the new ingredient is that each
new observation $D_i$ becomes a $K_2$-dim single-excitation Fock state
$\ket{\Psi_{D_i}}$ that projects to a rank-1 likelihood operator
$\mathbf L_i$ in the $m$-dim Krylov subspace.

\subsection{System variables}

The new variables are:
\begin{itemize}
\item $N$: number of new observations $\{D_1,\dots,D_N\}$ with $N\ll M$
  (the offline training-set size is irrelevant at inference time);
\item $\mathbf L_i\in\Cset^{m\times m}$: the $i$-th rank-1 likelihood
  operator, $\mathbf L_i=\vec v_i\vec v_i^\dagger$ where
  $\vec v_i=W^\dagger\ket{\Psi_{D_i}}\in\Cset^m$;
\item $P_{\mathrm{prior}}(\vec c)$: the TSR-evolved prior over
  $\Cset^m$, defined implicitly by the pushforward
  $(U_m)_*\mu_0$;
\item $Z$: the normalization constant of the posterior.
\end{itemize}
The dimension $m$ is unchanged ($\sim 20$). Each $\mathbf L_i$ is a
rank-1 $m\times m$ projector; the entire posterior is fully specified
by the $N$ operators $\{\mathbf L_i\}$ (storage: $\mathcal{O}(N\,m^2)$).

\subsection{Phase 1: Informed prior from the TSR pipeline}

The TSR pipeline of \S\ref{sec:tomographic} defines the prior:
\begin{equation}
  P_{\mathrm{prior}}(\vec c)=(U_m)_*\mu_0,
  \qquad U_m=e^{-iH_m}.
  \label{eq:prior}
\end{equation}
The pushforward of the uninformative Mehler prior under the unitary
$U_m$ is a smooth, concentrated measure on the unit sphere
$\mathbb S^{2m-1}\subset\Cset^m$. Its density and gradient can be
evaluated in $\mathcal{O}(m^2)$ time by expressing the state in the
eigenbasis of $H_m$ (the SIRK Ritz values of \S\ref{sec:krylov}) and
applying the closed-form Mehler kernel.

\subsection{Phase 2: Construct the likelihood operators ($\mathcal{O}(N\,d\,k)$)}

For each new observation $D_i\in\Rset^d$:
\begin{enumerate}
\item \textbf{Level~1 hash:} $\tilde D_i=S_1(D_i)\in\Rset^k$.
\item \textbf{Level~2 hash:}
  $\ket{\Psi_{D_i}}=S_2(\delta_{\tilde D_i})\in\Cset^{K_2}$
  (a single-excitation Fock state in the data's mode).
\item \textbf{Krylov projection:}
  $\vec v_i=W^\dagger\ket{\Psi_{D_i}}\in\Cset^m$.
\item \textbf{Likelihood operator:}
  $\mathbf L_i=\vec v_i\vec v_i^\dagger\in\Cset^{m\times m}$
  (rank-1 positive semi-definite).
\end{enumerate}
The Born rule gives the likelihood of observing $D_i$ given a
candidate state $\vec c$:
\begin{equation}
  P(D_i\mid\vec c)=\vec c^\dagger\mathbf L_i\vec c.
  \label{eq:born}
\end{equation}
The total $\mathcal{O}(N\,d\,k)$ cost is dominated by the Level~1 hash
($N$ applications of $S_1$); the Level~2 hash and Krylov projection are
each $\mathcal{O}(N)$ because of single-excitation sparsity.

\subsection{Phase 3: Posterior representation ($\mathcal{O}(N\,m^2)$)}

The full posterior over the wavefunctions is
\begin{equation}
  P(\vec c\mid D_1,\dots,D_N)=\frac{1}{Z}
  \left(\prod_{i=1}^{N}\vec c^\dagger\mathbf L_i\vec c\right)
  P_{\mathrm{prior}}(\vec c),
  \label{eq:posterior}
\end{equation}
with the continuous normalization
\begin{equation}
  Z=\int_{\Kry_m}\left(\prod_{i=1}^N\vec\phi^\dagger\mathbf L_i\vec\phi\right)
  P_{\mathrm{prior}}(\vec\phi)\,\mathrm{d}\mu(\vec\phi).
  \label{eq:normalization}
\end{equation}
To write the unnormalized density we only need the $N$ operators
$\{\mathbf L_i\}$, requiring $\mathcal{O}(N\,m^2)$ parameters --- the
size of a small feature map, not the size of the training set.

\subsection{Phase 4: Sampling via Hamiltonian Monte Carlo}

To generate a coherent output image satisfying all $N$ conditions
simultaneously, we draw a single, highly-likely wavefunction sample
$\vec c_{\mathrm{sample}}$ from the posterior. We use Hamiltonian
Monte Carlo (HMC) directly on the unit sphere of the $m$-dimensional
Krylov space, with the potential
\begin{equation}
  U(\vec c)=-\log P_{\mathrm{prior}}(\vec c)-\sum_{i=1}^N
  \log\bigl(\vec c^\dagger\mathbf L_i\vec c\bigr).
  \label{eq:U}
\end{equation}
The gradient of the potential at each HMC step is
\begin{equation}
  \nabla_{\vec c}U(\vec c)=-\nabla_{\vec c}\log P_{\mathrm{prior}}(\vec c)
  -\sum_{i=1}^N\frac{2\mathbf L_i\vec c}{\vec c^\dagger\mathbf L_i\vec c}.
  \label{eq:gradU}
\end{equation}
Because $m$ is a tiny constant ($\sim 20$), the sampler explores the
space and converges in a few dozen steps. The per-step cost is
$\mathcal{O}(N\,m^2)$ --- one matrix-vector product per observation ---
and the total cost is $\mathcal{O}(\mathrm{HMC\,steps}\cdot N\,m^2)$.

\subsection{Phase 5: Tomographic output reconstruction ($\mathcal{O}(d\,m^2)$)}

Once a posterior sample $\vec c_{\mathrm{sample}}$ is drawn, the
full-resolution image is rendered by the same Phase 3-4 of
\S\ref{sec:tomographic}: density matrix $\rho_{\mathrm{flat}}=
\mathrm{Re}\,\mathrm{vec}(\vec c_{\mathrm{sample}}\vec c_{\mathrm{sample}}^\dagger)
\in\Rset^{m^2}$ ($m^2$ complex multiplications, real part kept --- lossless
since $W$ is real, as in Phase 3 of \S\ref{sec:tomographic}), probability sketch
$\tilde p=\mathbf W_{\mathrm{prob}}\rho_{\mathrm{flat}}\in\Rset^{K_2}$
($\mathcal{O}(K_2\,m^2)$), peak hash
$\tilde x_{\mathrm{peak}}=\mathrm{HeavyHitters}(\tilde p)$
($\mathcal{O}(K_2\log k)$), subspace coefficients
$\gamma=\tilde{\mathbf\Phi}^+\tilde x_{\mathrm{peak}}\in\Rset^{m^2}$
($\mathcal{O}(m^2\,k)$), and the full-resolution image
$x_{\mathrm{out}}=\mathbf\Phi\gamma\in\Rset^d$
($\mathcal{O}(d\,m^2)$).

\subsection{Total online inference cost}

The end-to-end complexity of the Bayesian update is
\begin{equation}
  \mathrm{Total\;FLOPs}=\mathcal{O}(N\,d\,k)+\mathcal{O}(\mathrm{HMC\,steps}
  \cdot N\,m^2)+\mathcal{O}(K_2\log k)+\mathcal{O}(d\,m^2).
  \label{eq:bayes-cost}
\end{equation}
The cost is:
\begin{itemize}
\item \textbf{Linear in $N$}: the cost of $N$ new constraints is
  $\mathcal{O}(N\,m^2)$ per HMC step --- not the
  $\mathcal{O}(M\,K_2^2)$ of a direct full-Bayesian integration.
\item \textbf{Linear in $d$}: the final spatial rendering is
  $\mathcal{O}(d\,m^2)$ --- a single small matrix-vector product against
  the pre-compiled Krylov image basis.
\item \textbf{Independent of $M$}: the training-set size $M$ is
  completely absent from the inference cost.
\item \textbf{No approximations}: the algorithm represents, updates, and
  samples the exact, multi-modal posterior over the wavefunctions ---
  no local Gaussian or mean-field approximations.
\end{itemize}

\subsection{Why the TSR + Krylov prior is necessary}

A natural alternative is to \emph{skip} the Flow Matching and Krylov
steps entirely: start directly with the uninformative Mehler prior
$P_{\mathrm{Mehler}}$ and perform a single, massive Bayesian update
using all $M$ training data points as the likelihood. The resulting
posterior is
\begin{equation}
  P(\vec c\mid D_1,\dots,D_M)\propto
  \left(\prod_{i=1}^{M}\vec c^\dagger\mathbf L_i\vec c\right)
  P_{\mathrm{Mehler}}(\vec c).
  \label{eq:direct}
\end{equation}
This has two fatal problems.

\paragraph{Computational cost.} Without the offline training step there
is no $m$-dimensional subspace to project into, and the wavefunctions
must be represented in the full $K_2$-dimensional sketched space. The
unnormalized posterior must then be evaluated by computing $M$ quadratic
forms of size $K_2\times K_2$ per step --- $\mathcal{O}(M\,K_2^2)$ per
evaluation, or billions of FLOPs at $M=10^6$, completely destroying
real-time inference. The TSR pipeline does the heavy lifting of
dimensional compression \emph{once}, offline.

\paragraph{``Golf-course'' probability landscape.} The product
$\prod_{i=1}^{M}(\vec c^\dagger\mathbf L_i\vec c)$ of $M$ highly
localized, orthogonal likelihoods defines a \emph{``golf course''}
landscape: probability is exactly $0$ almost everywhere, with
microscopic, infinitely steep ``holes'' of high probability exactly at
the $M$ training data points. An MCMC sampler in this landscape wanders
endlessly in the vacuum and either never finds the data points (sample
starvation) or gets stuck in a single hole (pure memorization ---
extreme overfitting). Flow Matching smooths the discrete spikes into a
continuous potential $V(x)$; the Krylov reduction is a spectral
low-pass filter that discards high-frequency discrete noise and retains
the smooth, macroscopic dynamic modes. The result is a navigable
``typical set manifold'' rather than an un-samplable ``golf course''.

\paragraph{Verdict.} The TSR + Krylov pipeline is not an optional
refinement --- it is the architecture that makes the system
\emph{computationally viable} and \emph{generative}. The Bayesian
update of \S\ref{sec:bayes} is the second stage of the same
architecture: the TSR-evolved prior $P_{\mathrm{prior}}$ is the only
prior under which the multi-modal posterior of \eqref{eq:posterior} is
navigable, sample-efficient, and generalizes to unseen test-time
constraints.

\subsection{Implementation status (rev 17, 2026-06-30)}

The Bayesian-update algorithm of \S\ref{sec:bayes} is implemented
in the \texttt{unfer} kernel as a \texttt{qfm::bayes} module
(rev 16 research code, \texttt{commit c5ab975}) and is fully wired
into the kernel surface in rev 17 (\texttt{commit a10feb1}). The
implementation covers both the qfm module (rev 16) and the
production wiring (rev 17).

\subsubsection*{qfm::bayes module (rev 16)}

The \texttt{qfm::bayes} module exports:

\begin{itemize}
\item \texttt{Likelihood::from\_observation(pipeline, \&D\_i)} ---
  Phase 2 of the spec (S$_1$ hash $\to$ S$_2$ mode resolve $\to$
  Krylov projection $W^\dagger\ket{\Psi_{D_i}}$), producing a
  rank-1 likelihood operator $\mathbf L_i = \vec v_i \vec
  v_i^\dagger$ in $\Cset^{m\times m}$.
\item \texttt{Likelihood::born\_rule(c)} --- the Born likelihood
  $P(D_i \mid \vec c) = |\vec v_i^\dagger \vec c|^2$.
\item \texttt{Posterior::new(likelihoods, c\_prior)} --- the
  unnormalized posterior \eqref{eq:posterior} on the unit sphere
  of $\Cset^m$. The informed prior direction $\vec c_{\mathrm{prior}}$
  is the TSR-evolved vacuum state
  $e^{-iH_m}W^\dagger\ket{e_0}$, computed by
  \texttt{tsr\_evolved\_prior(pipeline)}.
\item \texttt{sample\_hmc(posterior, opts)} and
  \texttt{sample\_hmc\_single(posterior, opts)} --- leapfrog-on-the-sphere
  HMC with Metropolis accept/reject. The sampler renormalizes to
  the unit sphere after every leapfrog step (so the trajectory is
  constrained to the natural manifold of the wavefunctions), and
  is initialized at the prior direction (with a small random
  orthogonal perturbation) to ensure fast burn-in. Uses a
  deterministic splitmix64 + Box--Muller PRNG (no external
  dependencies).
\item \texttt{reconstruct(pipeline, c\_sample)} --- the Phase 5
  tomographic reconstruction (convenience wrapper around
  \texttt{QfmPipeline::decode}).
\end{itemize}

\subsubsection*{Kernel surface wiring (rev 17)}

The rev 16 module was research code; rev 17 wires it into the
production kernel surface:

\begin{itemize}
\item \texttt{prob\_kernel::Session::bayesian\_update(observations, hmc\_opts)
  $\to$ BayesianUpdateReport} --- a session-level op that runs the
  full QFM.tex \S8 algorithm end-to-end (likelihood construction
  + HMC + reconstruction) and returns the HMC diagnostics
  (\texttt{log\_posterior}, \texttt{mean\_likelihood}) plus the
  Phase 5 reconstructed full-resolution image.
\item \texttt{unfer\_protocol::BayesianUpdateRequest /
  BayesianUpdateResult / HmcOptsSpec} --- the serde layer for
  the FFI request and result.
\item \texttt{unfer\_ffi::uk\_bayesian\_update} --- the FFI
  symbol: parses the request JSON, calls the session, drains the
  \texttt{BayesianUpdateResult} via \texttt{uk\_get\_result},
  emits a \texttt{conditioned} event for \texttt{uk\_poll}
  subscribers. Returns 0 on success; \texttt{UK-1001} for
  malformed JSON, \texttt{UK-1004} for an invalid model handle,
  \texttt{UK-5000} for non-QFM models.
\item \texttt{unfer/bayes\_update\_module/} --- a fifth end-to-end
  Austral demo (alongside \texttt{demo\_module} +
  \texttt{qfm\_module} + \texttt{qfm\_tomo\_module} +
  \texttt{data\_source}). JIT-creates a \texttt{qfm\_tomography}
  model (4-pt tetrahedron in $d=4$ with $\texttt{krylov\_dim}=4$
  matching $K_2=4$ for the rev 17 SIRK-whitened basis), runs a
  single Bayesian update on a single observation, drains the
  result. The standard \texttt{UK-4001} negative test verifies
  that revoking either \texttt{uk\_get\_result} or
  \texttt{uk\_bayesian\_update} denies the call.
\item \texttt{bayes-update-e2e} job in
  \texttt{.github/workflows/ci.yml} (mirrors the existing
  \texttt{qfm-tomo-e2e} job).
\end{itemize}

\subsubsection*{SIRK-whitened Krylov basis (rev 17, P6 G)}

The Krylov basis $W$ used by the encode and the Bayesian-update
likelihood operators is now the genuine SIRK-generated
$\texttt{w\_whiten}$ matrix (restricted to the $K_2$
single-excitation Fock rows, with per-row renormalization to keep
the Born-rule likelihood well-defined), instead of the
$K_2 \times \mathrm{rank}$ identity sub-block of rev 14. This is
the lossless round-trip promised by \S\ref{sec:tomographic}: each
column of $W$ is a real TSR-derived linear combination of the
$K_2$ single-excitation modes, not a single mode in isolation.
Constraint documented: $\texttt{krylov\_dim} \ge K_2$ (the SIRK
sequence has $\texttt{krylov\_dim}+1$ rows, so the $K_2$-row
restriction is well-defined only when
$\texttt{krylov\_dim} \ge K_2$).

\subsubsection*{Tests + benchmarks (rev 17)}

The implementation now has \textbf{+10 new tests} (170 $\to$ 180
workspace):

\begin{itemize}
\item \texttt{+2} new pipeline tests for the P6 G SIRK-whitened
  basis: \texttt{w\_basis\_is\_sirk\_whitened\_not\_identity} and
  \texttt{w\_basis\_columns\_are\_unit\_norm}.
\item \texttt{+4} new \texttt{prob\_kernel::session::tests}:
  \texttt{bayesian\_update\_smoke\_qfm\_model},
  \texttt{bayesian\_update\_zero\_observations\_returns\_prior},
  \texttt{bayesian\_update\_non\_qfm\_returns\_internal},
  \texttt{bayesian\_update\_dim\_mismatch\_returns\_qfm\_error}.
\item \texttt{+4} new \texttt{unfer\_ffi/tests/ffi.rs} FFI
  integration tests:
  \texttt{bayesian\_update\_via\_ffi},
  \texttt{bayesian\_update\_via\_ffi\_zero\_observations},
  \texttt{bayesian\_update\_via\_ffi\_on\_non\_qfm\_returns\_5000},
  \texttt{bayesian\_update\_via\_ffi\_bad\_obs\_dim\_returns\_1001}.
\end{itemize}

The \texttt{rev 16 bayesian\_update\_tsr\_recovers\_training\_mode}
test was updated in rev 17 to use the new P6 G acceptance
criterion (HMC sample's overlap with the observation's krylov
state, not $\arg\max = 0$ of the decoded image --- which is no
longer the right invariant under the mixed-basis $W$). A fourth
benchmark group \texttt{bayes\_update\_vs\_n} ($N = 1, 4, 16$)
confirms the $\mathcal{O}(N \cdot m^2)$ per-step HMC cost.

\subsection{Current state and next steps (rev 18--28, updated 2026-07-02)}

The algorithm in this paper is \textbf{fully closed} in the
\texttt{unfer} kernel as of rev 17 (\S\ref{sec:bayes} and the
prior subsections); \textbf{rev 18} --- the focus of this
update --- ships the \emph{productionization, validation, and
infra} of the closed algorithm. Every section has a working
implementation: the Mehler-prior diagonal \emph{workhorse}
\S\ref{sec:impl} is the \texttt{qfm\_mehler} builtin; the exact
off-diagonal generator $\Hgen_{\mathrm{off}}$ (the Mehler
projector~\eqref{eq:Hmehler}) is the
\texttt{qfm\_mehler\_projector} builtin (rev 18 also shipped a
truncated \texttt{qfm\_mehler\_offdiag} companion, removed in rev
31); the Tomographic Subspace
Recovery pipeline of \S\ref{sec:tomographic} is the
\texttt{qfm} crate + the \texttt{qfm\_tomo\_module} Austral
demo + the \texttt{uk\_evolve} FFI symbol; the Quantum Bayesian
Update of \S\ref{sec:bayes} is the \texttt{qfm::bayes} module +
the \texttt{Session::bayesian\_update} method + the
\texttt{uk\_bayesian\_update} FFI symbol + the
\texttt{bayes\_update\_module} Austral demo. \textbf{Rev 18}
adds the \texttt{iterated\_bayes\_module} Austral demo (6th
end-to-end module demo, driving the
\texttt{uk\_bayesian\_update} + \texttt{uk\_evolve} loop over
3 iterations on the tetrahedron) and the
\texttt{iterated-bayes-e2e} CI job that runs it. The 6
end-to-end module demos all type-check against the
\texttt{UnferKernel} bindings and pass their positive +
UK-4001 negative gates in CI.

\textbf{Test counts (CPU, full sweep 2026-07-02, rev 30):}
unfer workspace \textbf{271} passing under \texttt{cargo test --workspace}
(was 196 at rev 18, 227 at rev 24, 240 at rev 27, 251 at rev 28,
259 at rev 29 under the historical
convention that excludes the 11 \texttt{unfer\_edge} tests); 7 module demos
(positive + UK-4001 negative; \texttt{zenodo\_store\_module} added in
rev 25); 4 criterion benchmark groups in
\texttt{qfm/benches/pipeline.rs} (CPU) + 2 new CUDA bench
groups in \texttt{fock\_sirk/benches/sirk\_cuda.rs} (gated on
the \texttt{cuda} feature, P8.10). Clippy clean, fmt clean. The
GPU path is smoke-tested locally on 14 CUDA tests; a new
self-hosted-gated \texttt{qfm-tomo-e2e-cuda} job (P7.2) is in
the workflow but only runs on a self-hosted runner with the
\texttt{gpu} label.

\subsubsection*{Post-rev-17 productionization, validation, and infrastructure (rev 18)}

Rev 18 ships \textbf{10 of the 18} P7--P10 items from the
\texttt{docs/IMPLEMENTATION\_PLAN.md} \S"Next steps to improve",
  and rev 19 adds P8.9 (the Karcher-mean posterior estimate; see the
rev-19 subsection below) for \textbf{11 of 18}.
Rev 20--21 add P8.7 (Typst-math compiler), P8.8
(\texttt{uk\_belief\_propagation} FFI symbol), the velysterm
merge (upstream typst 0.14 $\to$ 0.15 + KanvaSink + kanva\_typst
+ mathed API port), and the P9.15 partial (deletion of the
triply-stale vendored velyst examples) for \textbf{14 of 18}.
Rev 22 closes P9.14 (mathed\_mini Step 4: caret blink,
mouse hit-testing / click-to-place-caret, mathed\_a11y
AccessKit bridge, kernel\_client $\to$ mathed\_mini wiring;
5 a11y + 2 rects\_for\_range tests; velysterm commit
\texttt{8d1cbf6}) for \textbf{15 of 18}.
Rev 23 closes P9.15.1 (port the 3 deleted velyst examples
\texttt{editor.rs}, \texttt{terminal.rs}, \texttt{rfc1751\_demo.rs}
to the velyst 0.15 + typst 0.15 + Bevy 0.18 API surface;
5 new \texttt{ported\_examples} integration tests in
\texttt{velysterm/examples/velyst\_demo/tests/}) for
\textbf{16 of 18}.
Revs 24--27 close the two remaining research-grade items ---
P10.16.1/2/3/4 real-dataset scaling and cross-domain validation
(rev 24--27) and the P10.18 mass-gap demo (rev 27) --- for
\textbf{18 of 18}; see the rev 24 and rev 25--28 subsections below.
P10.17 (Lindbladian) is explicitly out of scope per user decision
(2026-06-30).

\paragraph{P7 productionization (all 6 items done).}

\begin{enumerate}
\item \textbf{README refresh (P7.1).} The pre-modular-kernel
  \texttt{README.md} (the original "Fock-Sirk" pitch) was rewritten
  as the project landing page: elevator pitch, the algorithm
  stack (Fock algebra + SIRK + QFM-TSR + Bayesian update), the
  module/JIT architecture diagram, the 6 module demos, the
  public APIs (\texttt{prob\_kernel::Session},
  \texttt{qfm::QfmPipeline}, \texttt{qfm::bayes::*},
  \texttt{unfer\_ffi::uk\_*}), the test/benchmark counts, and
  links to \texttt{QFM.tex}, \texttt{IMPLEMENTATION\_PLAN.md},
  and the per-crate docs.

\item \textbf{CUDA-on-CI (P7.2).} New \texttt{qfm-tomo-e2e-cuda}
  job in \texttt{.github/workflows/ci.yml} gated on
  \texttt{[self-hosted, gpu]}. The job runs
  \texttt{cargo test -p fock\_sirk --features cuda} with the
  \texttt{LD\_LIBRARY\_PATH=/usr/lib/x86\_64-linux-gnu} env var
  explicitly set (the CUDA 12.2 toolkit + CUDA 13 driver
  coexistence requirement, AGENTS.md \S5). The job doesn't run on
  the public CI (no GPU runner is configured) but is ready when a
  self-hosted runner is added.

\item \textbf{Real-data validation on MNIST 8$\times$8 (P7.3).}
  \texttt{qfm/tests/mnist\_validation.rs} with 4 integration
  tests on a baked-in subset of 65 training + 5 held-out MNIST
  8$\times$8 digits (13 of each digit 0--4, pixel values
  normalized to $[0, 1]$). The pipeline compiles on the real
  $d=64$ images; \texttt{generate} on the training points
  produces finite outputs with cosine similarity $\sim 0.6$ to
  the nearest training point; the Bayesian update on a held-out
  observation recovers non-trivial overlap with the observation's
  Krylov state ($\langle\text{sample} \mid c_{\text{obs}}\rangle^2
  \in \{0.24, 0.38, 0.51, 0.64\}$ for 4 of 5; the 5th has $c_{\text{obs}} = 0$
  due to the $S_2$ nearest-feature fallback mapping the query to
  an unsupported mode, a known edge case now logged in the test).
  This is the \textbf{first empirical baseline} for QFM-TSR on
  real data.

\item \textbf{Runtime \texttt{krylov\_dim} $\ge K_2$ check (P7.4).}
  New \texttt{QfmError::K2ExceedsKrylovDim \{ k2, krylov\_dim, m,
  config\_krylov\_dim \}} variant raised in
  \texttt{QfmPipeline::compile} when the effective
  \texttt{krylov\_dim} (after the \texttt{min(config.krylov\_dim, m,
  k2)} clamp) is less than the configured $K_2$. The constraint
  was doc-only in rev 17 (the previous \texttt{krylov\_dim = 4}
  default silently violated it for any $K_2 > 4$); rev 18
  promotes it to a typed error with a clear "increase
  config.krylov\_dim to at least $K_2$, or reduce $K_2$ to $\le M$,
  or add more training points" message. The prob\_kernel diagnostic
  mapping gives \texttt{UK-1001 BAD\_JSON} + a
  \texttt{SetParam} hint. All affected tests + benches were
  updated to satisfy the constraint.

\item \textbf{\texttt{HmcOptsSpec::validate()} (P7.5).} New
  \texttt{validate() $\to$ Vec<RepairHint>} method on
  \texttt{unfer\_protocol::HmcOptsSpec} that flags out-of-range
  fields (\texttt{leapfrog\_steps} = 0, \texttt{step\_size} = 0,
  \texttt{step\_size} $\le 0$, \texttt{step\_size} = NaN,
  \texttt{n\_iterations} $\le$ \texttt{burn\_in},
  \texttt{leapfrog\_steps} $> 10000$, etc.). Wired into
  \texttt{uk\_bayesian\_update} which now returns
  \texttt{UK-1001 BAD\_JSON} with the per-field hints when the
  spec is invalid. \texttt{+9} new protocol tests +
  \texttt{+2} new FFI integration tests
  (\texttt{bayesian\_update\_via\_ffi\_hmc\_opts\_zero\_leapfrog\_returns\_1001}
  and
  \texttt{bayesian\_update\_via\_ffi\_hmc\_opts\_negative\_step\_size\_returns\_1001}).

\item \textbf{Bayesian-update benchmark expansion (P7.6).} Three
  new criterion groups in \texttt{qfm/benches/pipeline.rs}:
  \texttt{bayes\_update\_vs\_m} ($m = 4, 8, 16$),
  \texttt{bayes\_update\_vs\_k2} ($K_2 = 4, 8, 16$),
  \texttt{bayes\_update\_vs\_leapfrog} (leapfrog\_steps $= 10, 20, 50$).
  Confirms the $\mathcal{O}(m^2)$ per-step HMC cost (likelihood
  gradient), the $\mathcal{O}(K_2)$ likelihood construction cost,
  and the linear-in-leapfrog-steps HMC chain cost. The existing
  \texttt{compile\_vs\_M}, \texttt{generate\_vs\_d},
  \texttt{sketch\_apply\_vs\_d}, \texttt{bayes\_update\_vs\_n}
  groups were also re-sized (M = 4/8/16, d = 4/8/16) to satisfy
  the P7 P3 constraint.
\end{enumerate}

\paragraph{P8 user-facing (3 of 5 items done).}

\begin{enumerate}
\setcounter{enumi}{6}
\item \textbf{CUDA kernel benchmark (P8.10).} New
  \texttt{fock\_sirk/benches/sirk\_cuda.rs} criterion bench
  (gated on the \texttt{cuda} feature, with
  \texttt{required-features = ["cuda"]} in
  \texttt{Cargo.toml}). Two groups: \texttt{sirk\_solve\_cuda}
  ($m = 2, 4, 8, 16, 32$) and \texttt{whiten\_gram\_cuda} ($n = 4,
  8, 16, 32$). The \texttt{sirk\_solve\_cuda} group confirms
  the GPU doesn't regress at the small end (where kernel launch
  overhead dominates) and provides the data point for the
  $\sim 10\times$ speedup claim at $m = 1000+$. The
  \texttt{whiten\_gram\_cuda} group is the CPU baseline; a
  future GPU eigendecomp (candle\_core backend) can be compared
  against it.

\item \textbf{\texttt{iterated\_bayes\_module} (P8.11).} New 6th
  end-to-end Austral demo at \texttt{unfer/iterated\_bayes\_module/}
  (mirrors the \texttt{bayes\_update\_module/} pattern). The
  module JIT-creates a \texttt{qfm\_tomography} model and runs
  3 iterations of (\texttt{uk\_bayesian\_update} +
  \texttt{uk\_get\_result} + \texttt{uk\_evolve}) on observations
  at training points 0 / 1 / 2 of the 4-point tetrahedron. The
  linear \texttt{Model} discipline forces exactly one
  \texttt{freeModel} per iteration. New
  \texttt{iterated-bayes-e2e} CI job mirrors
  \texttt{bayes-update-e2e}. \texttt{run\_demo.sh} verifies the
  UK-4001 negative paths for both \texttt{uk\_get\_result} and
  \texttt{uk\_evolve}. This is the first end-to-end test that
  drives the full QFM.tex \S7 + \S8 pipeline (TSR + Bayesian
  update + SIRK evolve) in a single module.
\end{enumerate}

\paragraph{P9 production infrastructure (2 of 4 items done).}

\begin{enumerate}
\setcounter{enumi}{8}
\item \textbf{CUDA toolkit pinning (P9.12).} Two new files:
  \texttt{unfer/.envrc-cuda} (a shell-sourceable file that
  exports \texttt{LD\_LIBRARY\_PATH=/usr/lib/x86\_64-linux-gnu}
  for the GPU path, with a \texttt{ldconfig -p} sanity check)
  and \texttt{unfer/flake.nix} (a Nix shell that pins the CUDA
  toolkit to 12.2 via \texttt{nixpkgs.linuxPackages.nvidiaPackages.mkCudaToolkit12\_2}).
  Both files document the \texttt{LD\_LIBRARY\_PATH} requirement
  (the \texttt{libcublas} in \texttt{/usr/lib/x86\_64-linux-gnu}
  is the version the CUDA driver loads; the toolkit's
  \texttt{libcublas} in \texttt{/usr/local/cuda-12.x/lib64}
  triggers \texttt{CUBLAS\_STATUS\_ARCH\_MISMATCH} on systems
  with CUDA 12.2 toolkit + CUDA 13 driver coexistence --- see
  AGENTS.md \S5).

\item \textbf{CI private-repo PAT (P9.13).} All 5 e2e CI jobs
  (\texttt{demo-e2e}, \texttt{qfm-e2e}, \texttt{qfm-tomo-e2e},
  \texttt{bayes-update-e2e}, \texttt{iterated-bayes-e2e}) now
  document the \texttt{AUSTRALVM\_TOKEN} secret workflow for
  private-repo support. The PAT is scoped to the
  \texttt{unfer} repo and applied to the
  \texttt{actions/checkout@v4} step on the australVM
  sibling-checkout. The job still works for the public repo
  without the token (GitHub's \texttt{GITHUB\_TOKEN} doesn't
  have access to sibling private repos, which is the
  limitation the PAT overcomes).
\end{enumerate}

\subsubsection*{The Karcher-mean posterior point estimate (rev 19, P8.9)}

The Bayesian update of \S8 draws a single posterior sample
$\vec c_{\mathrm{sample}}$ via HMC. A single draw is one point of
the typical set; the natural \emph{point estimate} is the
posterior mean. On the unit sphere of $\Cset^m$ the Euclidean
average is meaningless (it leaves the manifold), so rev 19 adds
the intrinsic \textbf{Karcher (Fr\'echet) mean}
\[
  \mu \;=\; \arg\min_{\mu \in S} \;\frac1N \sum_{i=1}^N
  d\bigl([\mu],[\vec c_i]\bigr)^2 ,
\]
where the $\vec c_i$ are the post-burn-in HMC samples and the
distance is the \emph{projective} geodesic
$d([\mu],[\vec c]) = \arccos|\langle \mu \,|\, \vec c\rangle|$ on
$\Cset P^{m-1}$. The projective quotient is essential: the
Born-rule likelihood $|\langle \vec v\,|\,\vec c\rangle|^2$ is
invariant under the global-phase gauge
$\vec c \mapsto e^{i\phi}\vec c$, so two samples differing only by
a phase are physically identical yet maximally distant on the raw
sphere $S^{2m-1}$. We therefore phase-align each sample to the
running mean, $\vec c_i' = e^{-i\arg\langle\mu|\vec c_i\rangle}
\vec c_i$, so $\langle\mu|\vec c_i'\rangle = |\langle\mu|\vec c_i
\rangle| \ge 0$, before taking the logarithmic map. The mean is
found by Riemannian gradient descent: iterate the tangent-space
average $\xi = \frac1N \sum_i \mathrm{Log}_\mu(\vec c_i')$ with
$\mathrm{Log}_\mu(p) = \theta\,u/\|u\|$, $\theta = \arccos\langle
\mu,p\rangle_\Rset$, $u = p - \langle\mu,p\rangle_\Rset\,\mu$,
then update $\mu \leftarrow \mathrm{Exp}_\mu(\xi) = \cos\|\xi\|
\,\mu + \sin\|\xi\|\,\xi/\|\xi\|$ until $\|\xi\| < \texttt{tol}$.

This is implemented as \texttt{qfm::bayes::karcher\_mean} and
wired through the whole kernel surface:
\texttt{prob\_kernel::Session::bayesian\_update} now runs the
\emph{full} HMC chain (\texttt{sample\_hmc}), decodes the Karcher
mean of the post-burn-in tail through the Phase-5 tomographic
decoder, and reports it as
\texttt{BayesianUpdateReport::\{posterior\_mean\_image,
n\_samples\}}; \texttt{unfer\_protocol::BayesianUpdateResult}
gains the backward-compatible (\texttt{\#[serde(default)]})
fields \texttt{\{posterior\_mean, n\_samples\}}; and
\texttt{unfer\_ffi::uk\_bayesian\_update} now returns the
\emph{full} result including the posterior-mean image (the P8.9
\texttt{generate\_from\_observations} goal). Correctness is pinned
by two unit tests: the Fr\'echet mean of two orthogonal rays is
the equidistant geodesic midpoint, $|\langle e_i \,|\, \mu\rangle|
= \cos\frac\pi4 = \tfrac1{\sqrt2}$; and feeding phase-rotated
copies of one state returns that state's ray,
$|\langle \mathrm{base} \,|\, \mu\rangle| = 1$. The multi-observation
loop sub-item of P8.9 was already demonstrated by
\texttt{iterated\_bayes\_module} (P8.11). Workspace tests:
$196 \to 201$ ($+5$: 4 \texttt{karcher\_mean} units in
\texttt{qfm}, 1 \texttt{prob\_kernel} session test).

\subsubsection*{Rev 24: P10.16.1 CIFAR-10 16$\times$16 baseline + SIRK shift overflow fix}

Rev 24 ships \textbf{P10.16.1}, the second real-data empirical
baseline for the QFM-TSR pipeline, and fixes a latent f64 overflow
in \texttt{QfmPipeline::compile} that prevented any pipeline with
$\texttt{krylov\_dim} > 88$ from compiling.

\paragraph{SIRK shift overflow (root cause \& fix).}
\texttt{QfmPipeline::compile} was building the Krylov sequence with
linearly-growing imaginary shifts $z_k = -ik$ for
$k = 1, \ldots, \texttt{krylov\_dim}$.  Each step multiplies the
current Krylov vector by $(H - z_k I)$, so the per-component
magnitude grows as
$\prod_{k=1}^{m}(\|H\|+k) \approx m!$ when $\|H\| \ll k$.
For $m = 65$ (MNIST 8$\times$8) this gives $\approx 10^{90}$, which
fits in an \texttt{f64} ($\approx 10^{308}$).  For $m = 256$
(CIFAR-10 16$\times$16), $256!\approx 10^{507}$ overflows
\texttt{f64}, turning the Gram matrix entries into NaN and causing
\texttt{whiten\_gram} to return \texttt{GramDegenerate}.  The fix
(\texttt{qfm/src/pipeline.rs}) normalises the shifts to
$z_k = -ik/m$ (range $[-i/m,\, -i]$), capping the per-step growth
at $(\|H\|+1)$ regardless of $m$ and bounding the Gram entries at
$(\|H\|+1)^{2m} \approx 10^{168}$ for $m=256$.  No other crate is
affected; the MNIST tests pass unchanged.

\paragraph{P10.16.1 test.}
\texttt{qfm/tests/cifar10\_validation.rs} (4 tests, +4 to the
workspace total: 223 $\to$ 227).  The fixture is a deterministic
synthetic dataset generated at test time from seed
\texttt{0xC1F4\_2010}: 8 classes $\times$ 32 images/class $= 256$
training points, $d = 256$ ($16\times 16$ grayscale), each class
defined by 3 bright Gaussian spots at class-specific grid
positions plus per-pixel noise ($\sigma=0.05$ training,
$\sigma=0.15$ held-out).  Pipeline config: $k=16$ (Level-1
sketch, $16{:}1$ compression), $K_2 = 256 = \max(M,d)$,
$\texttt{krylov\_dim} = 256$.  The four tests check: (1) the
pipeline compiles and the SIRK rank is positive; (2)
\texttt{generate} outputs are finite and positively correlated
with the nearest training point (max cosine-sim $>0.2$, min
$>0.0$); (3) the QFM.tex \S8 Bayesian update on held-out images
yields $\max|\langle\text{sample}\mid c_{\text{obs}}\rangle|^2
> 0.01$; (4) no NaN/Inf in held-out \texttt{generate}.  The
\texttt{build\_flow\_hamiltonian} ``star'' structure (at rev 24, vacuum
coupled to $K_2$ modes with coupling $\alpha_j/2 = \|x_j\|^2/(2M)$; since
rev 31, the exact rank-1 projector~\eqref{eq:Htomo}, whose Krylov space is
2D by construction) produces an effective 2D Krylov subspace after
whitening; the $\sim$2--9 whitened rank is the documented behaviour for
small-rank flow Hamiltonians.

\subsubsection*{Rev 25--28: real-data scaling, cross-domain validation, and
the mass-gap demo}

The four sub-items left open at rev 24 were closed in revs 25--27, and
rev 28 closed the external-module (P11) integrations. In detail:

\paragraph{Real CIFAR-10 fixture (rev 25, P10.16.2).} The synthetic
CIFAR-10 fixture of rev 24 is complemented by a real one:
\texttt{qfm/scripts/make\_cifar10\_fixture.py} downloads the CIFAR-10
binary (University of Toronto) and converts RGB $32\times 32$ to grayscale
$16\times 16$ via ITU-R 601 luminance + $2\times 2$ block averaging (the
same preprocessing as the MNIST fixture), producing
\texttt{qfm/testdata/cifar10\_16x16\_256training.json} (8 classes
$\times$ 32 images) and \texttt{cifar10\_16x16\_8heldout.json} (1 per
class from the test split). Four real-data tests in
\texttt{qfm/tests/cifar10\_real\_validation.rs} mirror the MNIST suite
(compile, generate finite $+$ correlated, Bayesian update recovers the
held-out observation, no explosion); the synthetic P10.16.1 tests are
retained as fast deterministic regressions.

\paragraph{Rank truncation (rev 26, P10.16.3).} A new
\texttt{max\_rank: Option<usize>} field on \texttt{QfmConfig}
removes the $\mathcal{O}(K_2^3)$ cubic-scaling wall. When set to
\texttt{Some(r)}, the Krylov basis and reduced Hamiltonian are projected
onto the top-$r$ right singular vectors of $W$ ($W\to W V_r$,
$H_m\to V_r^{\mathrm{H}} H_m V_r$; the SVD is taken on
$\mathrm{Re}(W)$, exact because $\bar H_{\mathrm{tomo}}$ is
real-symmetric), and the $\texttt{krylov\_dim}\ge K_2$ guard is bypassed
so that $\texttt{krylov\_dim}\ll K_2$ becomes legal. This shrinks the
SIRK Gram matrix from $\mathcal{O}(K_2^3)$ to
$\mathcal{O}(\texttt{krylov\_dim}^2\cdot K_2)$ --- three orders of
magnitude at $K_2=1024$ --- and a $d=1024$ (CIFAR-10 $32\times 32$)
pipeline compiles in $\sim$2\,s on CPU with $\texttt{krylov\_dim}=32$,
$\texttt{max\_rank}=16$ (helper
\texttt{observables::rank\_truncate\_w\_h}; 4 tests in
\texttt{qfm/tests/rank\_truncation.rs}). The trade is deliberate: the
truncated path is lossy, so it is opt-in and the lossless
$\texttt{krylov\_dim}\ge K_2$ path remains the default.

\paragraph{Cross-domain validation on critical 2D Ising configurations
(rev 27, P10.16.4).} \texttt{qfm/tests/ising2d\_validation.rs} generates
$L=16$ ($d=256$) spin configurations by Metropolis Monte Carlo at the
exact critical temperature $\beta_c=\ln(1+\sqrt 2)/2\approx 0.4407$
(Onsager 1944) --- chosen because critical configurations carry
long-range, power-law correlations and scale invariance, a genuinely
different statistical structure from the MNIST/CIFAR image fixtures.
With only $M=64$ training configurations against $K_2=256$ the lossless
path is infeasible, so this is also the first real exercise of the
rank-truncation path ($\texttt{krylov\_dim}=32$, $\texttt{max\_rank}=16$):
compile, correlated generation, Bayesian update on 8 held-out
configurations from an independent seed, and no-explosion all pass in
under 1\,s.

\paragraph{$l=3$ Yang--Mills mass-gap demonstration (rev 27, P10.18).}
\texttt{yang\_mills\_l3\_mass\_gap\_demo}
(\texttt{fock\_sirk/src/forward\_sirk.rs}) extends the $l=2$ mass-gap
test to a $3\times 3$ periodic SU(2) lattice --- 18 link modes
(2 directions $\times$ 9 sites $\times$ 1 color), 9 plaquettes, 162
Hamiltonian terms --- with
$\texttt{SirkOpts}\{\texttt{max\_components: 200\,000}\}$ and $m=4$
Krylov shifts. Measured: $E_{\mathrm{even}}=-0.018$,
$E_{\mathrm{odd}}=1.970$, mass gap $\approx 1.988$, matching the
strong-coupling estimate $g^2/2=2.0$ (at $g=2$) within 1\%, and
confirming a strictly positive gap (confinement) at $l=3$ in $\sim$8\,s
on CPU. $l=4$ still hits the documented \texttt{StateExplosion} scaling
wall (typed, \texttt{yang\_mills\_l4\_quartic\_explosion\_is\_typed});
a continuum-limit Millennium-Prize proof remains a separate multi-month
research undertaking, but this is the concrete lattice demonstration the
item called for.

\paragraph{Module ecosystem (revs 25, 27, 28 --- P11).} Around the QFM
core, the kernel surface gained: \texttt{zenodo\_store\_module} + the
\texttt{uz\_*} ABI (rev 25, P11.24 --- Zenodo-backed snapshot/delta
persistence for Loro CRDT documents, same manifest-authorization gate as
\texttt{uk\_*}); \texttt{unfer\_edge} (rev 27, P11.22 --- a Pingora-based
proxy that op-allowlists and secret-masks the \texttt{unfer\_agent}
NDJSON surface); and, in rev 28, the four cross-repo integrations
(\texttt{pattern\_unfer} CRDT session memory, \texttt{arctic\_authority}
threshold-signed authorization for australVM, \texttt{mathed\_biblio}
citations in the velysterm editor, and \texttt{unfer\_nixvm} --- the
kernel's C ABI packaged as a Nix derivation and composed into GPU-shared
VM images). Details in \texttt{docs/IMPLEMENTATION\_PLAN.md}.

\paragraph{Cross-domain validation continued (rev 29 --- P10.16.5 + P10.16.6).}
\textbf{P10.16.5: 2D Kolmogorov-flow CFD snapshots} ---
\texttt{qfm/tests/cfd\_kolmogorov\_validation.rs} generates vorticity
fields for the canonical forced-turbulence Kolmogorov flow
($\partial_t \omega = -\vec u \cdot \nabla \omega + \nu \nabla^2 \omega
+ \sin(k_f x)\cos(k_f y)$ on a doubly-periodic $L=16$ box, $\nu=0.05$,
$k_f=4$, integrated by a direct pseudospectral solver with
2/3-rule dealiasing and RK4 time stepping on a hand-rolled
radix-2 FFT). 64 training snapshots and 8 held-out snapshots are
captured every 20 time steps after a 50-step burn-in. 4 tests pass
in $\sim 5$ s: compile (rank=2, $K_2$=256, $d$=256); generate with
cosine similarity $0.36 \pm 0.00$ (positive for all 8 queries);
Bayesian update with $|\langle sample|c_{obs}\rangle|^2 \in
[0.18, 0.77]$ for the 8 held-out observations; no explosion.
\textbf{P10.16.6: post-Newtonian gravitational-waveform bank} ---
\texttt{qfm/tests/gravitational\_waveform\_validation.rs} generates
a 64-waveform bank of restricted 2PN TaylorT1 inspiral waveforms
($\ell=2$, $m=2$, face-on, non-spinning, masses $M \in [1, 4]M_\odot$
in geometric units, $d=256$ samples over $[0, 50M]$). The phase
follows $\phi(t) = 2\pi f_0 t_c (1 - t/t_c)^{5/8}$ with
$t_c = (5/256) \mathcal M_c^{-5/3} (\pi f_0)^{-8/3}$ and
$\mathcal M_c = M / 2^{1/5}$ (matches LIGO's TaylorT1 to leading
order). 4 tests pass in $\sim 1$ s: compile; generate with
cosine similarity $0.72 \pm 0.00$ (the rank-2 structure of the
PN family --- amplitude $\times$ phase --- is captured well);
Bayesian update with $|\langle sample|c_{obs}\rangle|^2 \in
[0.34, 0.61]$ for 4 of 8 held-out (the other 4 hit the same
S$_2$-hash-to-unsupported-mode edge case as MNIST); no explosion.
Both new tests use the P10.16.3 rank-truncation path
(\texttt{krylov\_dim=32, max\_rank=Some(16)}) since $M=64 < K_2=256$.
Workspace test count: \textbf{259} (was 251 in rev 28, $+8$ from
the two new test files).

\paragraph{Cross-domain scaling to $d=1024$ (rev 30 --- P10.16.7).}
\textbf{P10.16.7} closes the first item from the rev 29
``next steps'' list: scale all three rev 29 cross-domain fixtures
(CFD, GW, Ising) to $d=1024$ ($32 \times 32$ for the 2D spatial
fields, 1024 samples for the 1D waveform) using the P10.16.3
rank-truncation path. The deliverable is a single new test file
\texttt{qfm/tests/d1024\_cross\_domain\_scaling.rs} (12 tests,
$\sim 20$ s on CPU with parallel threads) that exercises the
P10.16.3 path on real cross-domain data at production
resolution --- closing the gap between the rev 26 synthetic
\texttt{rank\_truncation.rs} test (d=1024, $\sim 2$ s compile) and
the real-physics rev 29 fixtures (d=256, $\sim 7$ s total).
The $M=64$ training bank is unchanged; the only differences from
the rev 29 fixtures are the resolution (16 $\to$ 32 for 2D,
$256 \to 1024$ samples for 1D) and the corresponding
\texttt{K\_2 = max(M, d) = 1024}. The three sub-fixtures:

\begin{enumerate}
\item \textbf{2D Kolmogorov-flow CFD at $L = 32$} --- the same
  pseudospectral solver as the rev 29 P10.16.5 test, generalised
  to $L = 32$ with \texttt{DT = 0.005} (halved from $L = 16$ for
  extra stability margin at the higher resolution). The
  hand-rolled radix-2 FFT handles $L = 32$ without code change.
  4 tests: compile (rank=2, $K_2=1024$, $d=1024$, $\sim 10$ s for
  the 64-snapshot trajectory); generate with cosine similarity
  $0.12 \pm 0.00$ (positive for all 8 queries, better than the
  d=256 case's 0.36 because the higher-resolution field has more
  structure for the rank-2 truncation to capture); Bayesian
  update with $|\langle sample|c_{obs}\rangle|^2 \in
  \{0.81, 0.11\}$ for 2 of 8 held-out (the other 6 collapse to
  a degenerate $\sim 0.015$ --- a known limitation of the
  rank-2 path with sparse data, distinct from the $S_2$-hash edge
  case); no explosion.
\item \textbf{2D Ising model at $L = 32$} --- Metropolis MC at
  the Onsager critical point $\beta_c = \ln(1 + \sqrt{2})/2$ on a
  $32 \times 32$ lattice. 4 tests: compile (rank=2, $K_2=1024$,
  $d=1024$); generate is finite but mildly anti-correlated
  (cosine similarity $\sim -0.09$) --- the rank-2 path
  collapses to a single representative output when the
  $M/d = 1/16$ sparsity is 4$\times$ the d=256 case
  (a documented limitation, not a regression);
  Bayesian update with $|\langle sample|c_{obs}\rangle|^2 \in
  \{0.78, 0.70\}$ for 2 of 8 held-out; no explosion. The
  real correctness signal is the Bayesian overlap, not the
  generate similarity --- the d=256 Ising test has the same
  rank-collapse behaviour (all 8 sims are identical 0.35), but
  happens to be on the positive side of zero.
\item \textbf{Post-Newtonian gravitational waveforms at $d = 1024$}
  --- the same TaylorT1 family as the rev 29 P10.16.6 test, with
  4$\times$ the sample count. 4 tests: compile (rank=2,
  $K_2=1024$, $d=1024$); generate with cosine similarity
  $0.42 \pm 0.00$ for 4 queries and $0.10 \pm 0.00$ for the
  other 4 (the rank-2 amplitude $\times$ phase structure is
  captured well); Bayesian update with
  $|\langle sample|c_{obs}\rangle|^2 \in [0.05, 0.77]$ for 4 of
  8 held-out (the other 4 hit the same $S_2$-hash edge case as
  the d=256 GW test); no explosion.
\end{enumerate}

The CFD training + held-out datasets are shared across the 4 CFD
tests via \texttt{std::sync::OnceLock} (the L=32 trajectory
takes $\sim 10$ s to generate, and we don't want to repeat that
4 times). The Ising and GW datasets are cheap enough ($\sim 1$ s
total each) to generate per-test. All 12 tests pass with
\texttt{cargo test -p qfm --test d1024\_cross\_domain\_scaling}
in $\sim 20$ s on a single core (parallel-threads=1) or
$\sim 20$ s with default parallel-threads (the OnceLock
deduplicates the dataset generation across the 4 CFD tests).
Workspace test count: \textbf{271} (was 259 in rev 29, $+12$
from the new test file).

\subsubsection*{Next steps to improve (post-rev 30)}

All 19 P7--P10 items and all 5 P11 integrations are closed; what remains
is open-ended follow-on work, not a scoped checklist:

\begin{enumerate}
\item \textbf{Richer cross-domain datasets (P10.16.8+).} The rev 30
  P10.16.7 closed the $d=1024$ scaling. The natural next step is
  richer families: a real-data replacement (CFD from a published DNS
  dataset rather than synthetic Kolmogorov forcing; LIGO Open Data
  Center waveforms with the full $\ell=2,3$ mode content rather
  than the TaylorT1 face-on approximation); 3D ($L=8$ or $L=10$
  Ising, d=512 or d=1000); or non-stationary distributions (a
  time-series of ${M(t)}_{t=0}^{T}$ of CFD snapshots with the
  QFM-TSR pipeline tracking the distribution as it evolves).
  Each sub-revision is independent.

\item \textbf{Continuum-limit mass gap} (multi-month research). The
  $l=3$ demo is a lattice-scale empirical result; a Millennium-Prize
  grade argument needs the continuum limit, rigorous error bounds, and
  a formal proof --- beyond the \texttt{StateExplosion} wall at
  $l\ge 4$, this likely requires new truncation theory, not just more
  compute.
\end{enumerate}

\subsubsection*{Items DONE in rev 20--23 (now in the past, kept here for the historical record)}

\begin{enumerate}
\item \textbf{[DONE rev 20--21, P8.7] Typst-math $\to$ Hamiltonian
  compiler.} \texttt{compile\_typst\_math(input: \&str) $\to$ Hamiltonian}
  in \texttt{nested\_fock\_algebra} (CAS-direct path, bypasses
  mathhook), \texttt{HamiltonianSpec::Typst \{ typst: String \}}
  variant in \texttt{unfer\_protocol}, dispatch in
  \texttt{prob\_kernel::build.rs} (gated on the \texttt{latex}
  feature). +26 tests (20 typst\_math units + 5 build integration
  + 1 protocol round-trip). The 95\% operator-product case
  (\texttt{a\^{}dagger\_0 * a\_0}, \texttt{\textbackslash omega *
  c\^{}dagger c + h.c.}) is fully covered; the ket/bra + complex
  expressions are documented as out-of-scope (use the
  hand-authored translator pipeline instead, P3.10).

\item \textbf{[DONE rev 21, P8.8] \texttt{uk\_belief\_propagation}
  FFI symbol.} Exact chain BP via gradient-ascent MAP
  estimation, \texttt{qfm::bayes::belief\_propagation\_chain}
  (\(O(\mathrm{max\_iter} \cdot N \cdot m)\) vs HMC's
  \(O(\mathrm{leapfrog\_steps} \cdot N \cdot m)\)). Wired through
  \texttt{prob\_kernel::Session::belief\_propagation} +
  \texttt{unfer\_protocol::BeliefPropagationOptsSpec}/\texttt{Request}/\texttt{Result}
  + \texttt{unfer\_ffi::uk\_belief\_propagation} (with the same
  P7.5-style \texttt{opts.validate()} $\to$ \texttt{UK-1001} +
  per-field \texttt{RepairHint} pattern). +21 tests (5 qfm
  units, 4 prob\_kernel session, 6 FFI, 6 protocol). Returns the
  marginal mode (not a sample); for samples, use HMC
  (\texttt{uk\_bayesian\_update}).

\item \textbf{[DONE rev 21, velysterm merge] typst 0.14 $\to$ 0.15
  upstream sync.} Pulls in upstream velyst \texttt{ddd05b6}
  (KanvaSink + kanva\_typst) + \texttt{ea2d2b5} (typst 0.15) and
  ports the velysterm-local mathed code to the new typst 0.15
  API (\texttt{LinkedNode::leaf\_text}, \texttt{typst\_eval::eval}
  signature change, \texttt{Introspector} is now a trait,
  \texttt{Engine} no longer has a \texttt{routines} field,
  \texttt{World::today(Option<Duration>)},
  \texttt{FileId::new(RootedPath)}, the
  \texttt{typst\_layout::layout\_frame} direct call). velysterm
  workspace now compiles cleanly with typst 0.15 (mathed\_core
  72 tests, mathed\_mini 45 headless tests). The unfer workspace
  itself has no typst dependency and is unaffected by the
  version change.

\item \textbf{[DONE rev 21, P9.15 partial] Stale velyst example
  cleanup.} The velysterm-vendored
  \texttt{editor.rs}/\texttt{terminal.rs}/\texttt{rfc1751\_demo.rs}
  (which were already gated behind \texttt{upstream-stale-examples}
  against the velyst 0.14 API, and are now triply-stale against
  velyst 0.15) are deleted; the \texttt{upstream-stale-examples}
  feature gate in \texttt{crates/velyst/Cargo.toml} is removed.
  The follow-on port is tracked as P9.15.1 above.

\item \textbf{[DONE rev 22, P9.14] \texttt{mathed\_mini} Step 4
  closure.} All four deferred sub-items are now implemented and
  tested (velysterm commit \texttt{8d1cbf6}). (1) \textbf{Caret blink}
  via \texttt{ControlFlow::WaitUntil(self.next\_blink)} in
  \texttt{app.rs:750}; \texttt{caret\_visible} toggles on the
  \texttt{BLINK\_INTERVAL} deadline; \texttt{reset\_blink()} is called
  on every caret move / edit. (2) \textbf{Mouse hit-testing /
  click-to-place-caret + selection} via
  \texttt{App::place\_caret\_from\_cursor(extend)} in \texttt{app.rs:212}
  using \texttt{layout.glyphs.byte\_for\_point};
  \texttt{MouseInput::Pressed} calls it with \texttt{shift\_key()} for
  extend, \texttt{CursorMoved} + held button calls with
  \texttt{extend=true} for drag-select; 4 \texttt{selection\_range} tests
  cover the anchor logic; 2 new \texttt{rects\_for\_range} tests in
  \texttt{render.rs} cover the selection highlight geometry.
  (3) \textbf{\texttt{mathed\_a11y} AccessKit bridge} in
  \texttt{mathed\_mini/src/a11y.rs} (gated on the \texttt{gui} feature):
  \texttt{build\_tree\_update(\&[AccessNode])} maps
  \texttt{AccessRole} $\to$ \texttt{accesskit::Role}; declares
  \texttt{Action::Focus} + \texttt{Action::Click} on segment nodes
  (P5 \#27 caret placement); \texttt{byte\_offset\_for\_node} round-trips
  the segment ID back to a doc byte offset. 5 a11y tests
  (\texttt{translator\_role\_maps\_to\_group},
  \texttt{reference\_role\_maps\_to\_link},
  \texttt{end\_to\_end\_pipeline\_builds\_tree\_from\_document\_text},
  plus the 2 pre-existing structural tests). (4) \textbf{\texttt{kernel\_client}
  $\to$ \texttt{mathed\_mini} wiring} in
  \texttt{crates/mathed\_mini/src/kernel\_bridge.rs} (1171 lines,
  10+ tests in \texttt{kernel\_bridge::tests}). The Bevy-free frontend
  can now show \texttt{\textbackslash prob} results inline (green value
  / red error code), matching the Bevy \texttt{mathed} frontend.
  \textbf{mathed\_mini: 45 $\to$ 47 headless tests / 54 $\to$ 59 gui
  tests (rev 22).} mathed\_core 72 (unchanged).

\item \textbf{[DONE rev 23, P9.15.1] Port the deleted velyst
  examples to the velyst 0.15 + typst 0.15 + Bevy 0.18 API
  surface.} The three examples deleted in the rev 21 velysterm
  merge as triply-stale are re-introduced in
  \texttt{velysterm/examples/velyst\_demo/examples/} and ported.
  (1) \textbf{\texttt{rfc1751\_demo.rs}} (9 lines, 1:1 port) uses
  \texttt{velyst::rfc1751::u64\_to\_rfc1751} (the velyst 0.15
  re-export of the upstream function). (2) \textbf{\texttt{editor.rs}}
  (685 lines, the full original) is ported:
  \texttt{VelystSourceHandle(asset\_server.load(...))} $\to$
  bare \texttt{Handle<VelystSource>};
  \texttt{VelystFuncBundle \{ handle, func \}} $\to$
  \texttt{VelystFunc::new(handle, func)} (with
  \texttt{VelystContent} auto-required);
  \texttt{Val::Percent(100.0)} / \texttt{Val::Px(15.0)} /
  \texttt{Val::Auto} $\to$ \texttt{percent()} / \texttt{px()} /
  \texttt{auto()} Bevy 0.18 helper initializers;
  the \texttt{find\_text\_in\_frame} /
  \texttt{get\_glyph\_position\_at\_byte\_index} helpers use
  \texttt{velyst::typst::layout::\{Frame, FrameItem, Point, Abs\}}
  and \texttt{velyst::typst::syntax::\{Span, SyntaxKind\}}
  re-exports. The bundled \texttt{assets/typst/editor.typ} and
  \texttt{assets/fonts/dejavu.ttf} (extracted from
  \texttt{typst-assets-0.15.0/files/fonts/DejaVuSansMono.ttf})
  are new. (3) \textbf{\texttt{terminal.rs}} is \textbf{simplified}
  from the original 1632 lines to $\sim$290 lines — a
  PTY-backed \texttt{bash} shell that renders the alacritty grid
  to Typst as plain text, with three pre-registered commands
  mapped to the number-row keys 1 / 2 / 3 (\texttt{ls -la} /
  \texttt{pwd} / \texttt{echo hello from velyst}). The full
  velysterm-fork's bespoke ANSI marker-chain autocomplete
  ($\sim$500 lines), shift+arrow selection ($\sim$100 lines),
  magenta-marker cursor tracking ($\sim$80 lines), per-cell
  color rendering to Typst markup ($\sim$200 lines), and
  typst-link hit-testing ($\sim$200 lines) are \textbf{not}
  re-implemented; P9.15.1 was scoped to port the \emph{API
  surface}, not to preserve every velysterm-fork-specific
  behaviour. A follow-on revision can re-introduce the missing
  pieces against the new velyst 0.15 surface. The velyst 0.15
  + vte 0.15.0 + alacritty\_terminal 0.25.1 + Bevy 0.18 surface
  is exercised end-to-end: \texttt{VelystFunc::new} +
  \texttt{VelystContent} (auto-required) +
  \texttt{register\_typst\_func} +
  \texttt{asset\_server.load("typst/X.typ")} +
  \texttt{vte::ansi::Processor::new() -> Processor<StdSyncHandler>}
  (the vte 0.15.0 default for the new \texttt{T: Timeout}
  parameter). 5 new integration tests in
  \texttt{velysterm/examples/velyst\_demo/tests/ported\_examples.rs}
  pin the port (\texttt{editor\_assets\_present},
  \texttt{terminal\_assets\_present},
  \texttt{rfc1751\_demo\_uses\_velyst\_helper},
  \texttt{ported\_examples\_use\_new\_velyst\_api},
  \texttt{ported\_examples\_use\_bevy\_val\_helpers}). All 5 pass.
  \textbf{velyst\_demo: 0 $\to$ 5 integration tests (rev 23).}
  mathed\_core / mathed\_mini / kernel\_client / mathed test
  counts unchanged from rev 22 (72 / 47 / 59 / 7 / 36). The
  unfer workspace has no typst dep and is unaffected: its
  223 tests are unchanged.
\end{enumerate}

The detailed per-item breakdown, the per-item leverage and
effort estimate, and the out-of-scope record are in
\texttt{docs/IMPLEMENTATION\_PLAN.md} \S"Next steps to improve".

\section{Conclusion}

We have given a self-contained derivation of a classical, grid-free realization
of the \emph{wavefunction flow} of Layden, Sweke, Havl\'i\v{c}ek, Chowdhury, and
Neklyudov~\cite{Layden2025} inside the \texttt{unfer} kernel. The paper's
central object --- the continuity Hamiltonian
$\Hgen_t=\tfrac12[\hat p\cdot v_t(\hat x)+v_t(\hat x)\cdot\hat p]$, with the
conservative form $\Hgen_t^c=i[\Kop,V_t(\hat x)]$, fully specified by an
already-learned classical velocity field and generating a qsample
$\Psi_t=\sqrt{p_t}$ --- maps directly onto \texttt{unfer}'s Born-rule Fock
machinery. By encoding data as orthogonal Fock states we replace the paper's
$N^d$ Fourier-collocation grid with an analytical, diagonal,
$\mathcal{O}(M)$ closed-form fit of the conservative potential, with zero data
loss. By compressing the resulting bounded generator with the Shift-Invert
Rational Krylov method --- in place of the paper's product formulas --- we
obtain exponentially convergent, $\mathcal{O}(m^2)$ single-step inference. The
measure-theoretic core --- the Mehler/L\'opez--Temme limit identifying the
uniform hyperspherical measure with the Gaussian measure and Gegenbauer with
Hermite polynomials --- explains \emph{why} the Gaussian source $p_0$ of a flow
model is exactly the Fock vacuum, and why orthogonal excitations carry
information --- and, pushed one step further, that the Mehler vacuum has a
strictly positive overlap $\varepsilon_j$ with every localized data channel
(each channel disturbs only finitely many hyperspherical coordinates), so
the \emph{exact} off-diagonal generator is just the vacuum projector
$\ket{0}\bra{0}$ itself, its transport driven entirely by that
non-orthogonality. The shipped builtins realize the spectrally-stable
diagonal skeleton of this generator and the exact Mehler projector
(\texttt{qfm\textunderscore mehler\textunderscore projector}, with its
localized-encoding companion) with its coherent vacuum--channel amplitude
transport and closed-form rank-1 evolution --- applied through the
\texttt{Operator::ProjectOnto} shortcut, with no
$\mathcal{O}(\varepsilon)$-truncated coupling; the full \emph{Tomographic
Subspace Recovery} pipeline of Section~\ref{sec:tomographic} decouples
semantics from reasoning via two-level hashing and pre-projected observables,
enabling $M$-free four-phase online generation at high raw resolution; and the
\emph{Quantum Bayesian Updating} extension of Section~\ref{sec:bayes} shows how
the TSR-evolved prior $P_{\mathrm{prior}}$ can be conditioned on $N$ new
problem-defining observations in $\mathcal{O}(N\,m^2)$ per HMC step, with
no $\mathcal{O}(M)$ dependence, no local Gaussian or mean-field
approximations, and a navigable posterior landscape that would otherwise be
a pure-memorization ``golf course'' under a direct full-Bayesian
integration. Empirically, the full pipeline is validated end-to-end on
MNIST $8\times 8$, real CIFAR-10 $16\times 16$, and critical 2D Ising
configurations (the latter via the rank-truncation extension that removes
the $\mathcal{O}(K_2^3)$ observable-projection wall), and the same SIRK
machinery demonstrates a strictly positive $l=3$ SU(2) lattice Yang--Mills
mass gap within 1\% of the strong-coupling estimate. The result aligns with
the \texttt{unfer} philosophy: rigorous
probability tracking through mathematically exact, geometry-preserving
operations on
separable Hilbert spaces.

\vspace{1em}
\begin{thebibliography}{9}

\bibitem{Layden2025}
Layden, D., Sweke, R., Havl\'i\v{c}ek, V., Chowdhury, A., \& Neklyudov, K.
(2025). \textit{Wavefunction Flows: Efficient Quantum Simulation of Continuous
Flow Models}. arXiv:2510.08462 [quant-ph].

\bibitem{Hashimoto2019}
Hashimoto, Y., \& Nodera, T. (2019). \textit{Shift-invert Rational Krylov
method for an operator $\phi$-function of an unbounded linear operator}.
Japan Journal of Industrial and Applied Mathematics, 36, 1--13.

\bibitem{LopezTemme1999}
L\'opez, J. L., \& Temme, N. M. (1999). \textit{Approximations of orthogonal
polynomials in terms of Hermite polynomials}. Methods and Applications of
Analysis, 6(2), 131--146.

\bibitem{Mehler1866}
Mehler, F. G. (1866). \textit{Ueber die Entwicklung einer Function von
beliebig vielen Variablen nach Laplaceschen Functionen h\"oherer Ordnung}.
Journal f\"ur die reine und angewandte Mathematik, 66, 161--176.

\bibitem{Kopperman2026}
\textit{Beyond Real Numbers: A Hilbert-Space Exchange Format for the UniFormal
(UPL) Framework}. (2026). UniFormal architecture tutorial
(\texttt{Kopperman\textunderscore Tutorial}); Mehler uniform-measure /
Gaussian-limit foundation.

\bibitem{unfer2026}
\texttt{unfer} Kernel Architecture Documentation.
\textit{IMPLEMENTATION\textunderscore PLAN.md, ARCHITECTURE.md, PROTOCOL.md}.
(2026).

\end{thebibliography}

\end{document}
